\documentclass[11pt]{article}
\usepackage[T1]{fontenc}
\usepackage{lmodern}
\usepackage[margin=1in]{geometry}
\usepackage{amsmath,amssymb,amsthm,mathtools,microtype,booktabs,array}
\usepackage[dvipsnames]{xcolor}
\usepackage[colorlinks=true,linkcolor=MidnightBlue,citecolor=MidnightBlue,urlcolor=MidnightBlue]{hyperref}
\usepackage{fancyhdr}
\pagestyle{fancy}\fancyhf{}
\fancyhead[L]{\small Sharp stability and exact cap optima}
\fancyhead[R]{\small Rademacher sums}
\fancyfoot[C]{\thepage}
\setlength{\headheight}{14pt}
\setlength{\parskip}{2pt}
\setlength{\emergencystretch}{2em}
\allowdisplaybreaks
\newtheorem{theorem}{Theorem}[section]
\newtheorem{lemma}[theorem]{Lemma}
\newtheorem{proposition}[theorem]{Proposition}
\newtheorem{corollary}[theorem]{Corollary}
\theoremstyle{remark}\newtheorem{remark}[theorem]{Remark}
\newcommand{\Pp}{\mathbb P}
\newcommand{\E}{\mathbb E}
\newcommand{\R}{\mathbb R}
\newcommand{\code}[1]{\texttt{\detokenize{#1}}}
\hypersetup{pdftitle={Sharp stability and exact cap optima for Rademacher sums},pdfauthor={Dev D. Goyal}}
\title{\textbf{Sharp stability and exact cap optima\\for Rademacher sums}}
\author{Dev D. Goyal}
\date{1 November 2026}

\begin{document}
\maketitle
\begin{abstract}
For a normalized finite Rademacher sum, we study the central support gap
and normal-approximation discrepancy under an upper bound on every
coefficient. Our principal result gives sharp quantitative geometry for
large gaps: if $|b_i|\le a$, every signed sum has magnitude at least
$x>a/2$, and $C=\{i:|b_i|\ge x\}$, then $C$ has positive odd size $k$
and the squared Euclidean distance to the flat vector on $C$ is at most
$2-2\sqrt{1-(a-x)^2}$. Equality is exactly an odd saturated core with
one residual coefficient. Combining the finite-slice geometry with a
sharp tail-compression argument gives a uniform stability theorem for
the exact support-gap optimum at every cap. Its optimal linear constant
is $\sqrt{2+\sqrt2}$, attained by splitting a residual into two equal
parts. We determine the exact local nonlinear modulus on the saturated
branch and the exact fixed-core flat modulus, including all equality
cases and the admissibility restrictions imposed by the cap.

We extend the exact full Kolmogorov cap optimum from the previously
proved range $a\ge1/\sqrt3$ to every $a\ge1/\sqrt5$, with all
optimizing coefficient vectors, thresholds, and transitions. The two
new cap intervals alternate between a saturated core with one residual
and the next flat vector. The probability proofs combine two- and
four-coordinate localization with a four-shift lower-discrepancy
envelope, certified on $27972$ rational boxes by validated interval
arithmetic. They use explicitly stated Keller--Klein, Hollom--Portier,
and Tomaszewski inputs. The deterministic stability results are
independent of those external theorems and certificates. Lean 4 checks
the actual finite-vector inverse, all-cap stability and optimizer
selection, exact saturated modulus and equality, fixed-flat inequalities,
and concrete sharpness examples. Writing $c(t)=\Phi(t)-1/2$ and
$\mathfrak G(a)$ for the optimal gap, the lower-discrepancy optimum equals
$c(\mathfrak G(a))$ throughout the new range, so every lower optimizer
has an empty strict central interval. The probability optimization and its
interval certificates remain outside this kernel-checked scope.
\end{abstract}

\section{Problem, contributions, and prior work}
Let $\varepsilon_i$ be independent uniform signs and let
\[
 X_b=\sum_{i=1}^N b_i\varepsilon_i,\qquad
 \sum_{i=1}^N b_i^2=1,\qquad |b_i|\le a\le1.
\]
The number of coefficients is arbitrary and finite. A cap $a$ need not
be attained by a coefficient. Write
\[
 g(b)=\min_{\varepsilon\in\{-1,1\}^N}
       \left|\sum_i b_i\varepsilon_i\right|,
 \qquad
 \mathfrak G(a)=\sup_{N,b}g(b).
\]
Signs of coefficients may always be absorbed into the random signs.
Vectors of different lengths are compared after arranging their absolute
coefficients in decreasing order and padding with zeros. We write
$\operatorname{dist}(b,v)$ for the resulting Euclidean distance, and
occasionally abbreviate it by $\|b-v\|_2$ when this convention is understood.
Thus $(a^k,r)$ denotes $k$ copies of $a$ followed by one residual $r$,
with zero residuals omitted. The flat vector of size $k$ has all its
nonzero coefficients equal to $k^{-1/2}$.

The earlier research package~\cite{Prior} established the exact value
\begin{equation}\label{eq:intro-gap}
 \mathfrak G(a)=\max\left\{\frac1{\sqrt{n+2}},\,
 a-\sqrt{1-na^2}\right\},\qquad
 \frac1{\sqrt{n+2}}\le a\le\frac1{\sqrt n},\quad n\text{ odd}.
\end{equation}
It also classified the maximizers: the winning type is either the next
odd-flat vector or an odd saturated core with one residual, and both
types maximize at the unique transition. We include the complete short
proof in Section~\ref{gap:section}; this exact value is an established
starting point, not a new theorem of the present extension.

\subsection{New structural results}
The question addressed here is quantitative: how much coefficient
structure does a nearly maximal gap force? A variance estimate alone
would typically yield a square-root dependence on the extremal deficit.
The relevant majority slice instead has a sharp linear geometric bound.
Theorem~\ref{thm:sharp-gap-inverse} proves, for every certified gap
$x>a/2$ and the actual core $C=\{i:|b_i|\ge x\}$,
\[
 \left\|(|b_i|)_i-|C|^{-1/2}\mathbf1_C\right\|_2^2
 \le2-2\sqrt{1-(a-x)^2}.
\]
The cap can be unattained and the certified gap can be smaller than the
actual one. Equality determines the coefficients completely. The strict
threshold $1/2$ is optimal: the normalized vector $(2,1,1,1)/\sqrt7$
has gap equal to half its largest coefficient but has an even core.

Theorem~\ref{thm:all-cap-gap-stability} turns this geometry into stability
for the exact optimum~\eqref{eq:intro-gap}. If
$x=g(b)>\max\{a/2,(n+4)^{-1/2}\}$, the core has size $n$ or $n+2$,
and the appropriate candidate lies within
\[
 \sqrt{2+\sqrt2}\,\bigl(\mathfrak G(a)-g(b)\bigr)
\]
of $b$. The constant is best possible. Below the difference between the
two competing values, the candidate must be the actual unique optimizer;
at a transition both candidates are necessary. The fixed-flat curve in
Proposition~\ref{prop:fixed-flat-gap} and the exact saturated modulus in
Theorem~\ref{thm:exact-saturated-stability} further describe the equality
geometry. In particular, the optimal local coefficient on a fixed
saturated branch is $1$; the larger universal constant appears when the
entire residual splits into two equal coefficients.

\subsection{Exact normal discrepancy down to the five-flat cap}
Let $\Phi$ and $\varphi$ denote the standard normal distribution and
density, and put $c(x)=\Phi(x)-1/2$. For $x\ge0$ define
\[
 U_b(x)=\tfrac12\Pp(|X_b|\le x)-c(x).
\]
For $x>0$, define
$L_b(x)=c(x)-\tfrac12\Pp(|X_b|<x)$, and set $L_b(0)=0$.
The distinction between the closed and open central events is essential
at atoms. Symmetry and continuity of $\Phi$ give
\[
 d_K(X_b,N(0,1))=\max\{\sup_{x\ge0}U_b(x),\sup_{x\ge0}L_b(x)\}.
\]
Write $\mathfrak D(a)$, $\mathfrak L(a)$, and $\mathfrak K(a)$ for the
cap suprema of $U$, $L$, and the full Kolmogorov distance respectively.

The new probability results determine the full optimum on the next two
cap intervals. Combining them with the established large-cap formulas
gives the following explicit profile.

\begin{theorem}[The full cap optimum for $a\ge1/\sqrt5$]
\label{thm:full-cap-profile}
For $1/\sqrt5\le a\le1$,
\begin{equation}\label{eq:full-cap-profile}
\mathfrak K(a)=
\begin{cases}
\max\{c(1/\sqrt5),\,3/16-c(\sqrt{1-4a^2})\},
 &1/\sqrt5\le a\le1/2,\\[2pt]
\max\{3/16,\,c(a-\sqrt{1-3a^2})\},
 &1/2\le a\le1/\sqrt3,\\[2pt]
\max\{c(1/\sqrt3),\,1/4-c(\sqrt{1-2a^2})\},
 &1/\sqrt3\le a\le1/\sqrt2,\\[2pt]
\max\{1/4,\,c(a-\sqrt{1-a^2})\},
 &1/\sqrt2\le a\le1.
\end{cases}
\end{equation}
When the constant branch is strictly larger, the unique maximizing type
is the next flat vector, of size $5,4,3,2$ respectively. When the varying
branch is strictly larger, it is the saturated vector
$(a^n,\sqrt{1-na^2})$ with $n=4,3,2,1$ respectively. Each interval
has one transition; precisely both types maximize there.
Winning even cores attain their value in $U$ at
the residual threshold, odd cores in $L$ at the support gap. A flat
vector attains its value in $U$ at zero when its size is even, and
in $L$ at its support gap when its size is odd.
\end{theorem}

The first two rows and their complete equality statements are new:
they are proved in Theorems~\ref{thm:full-five-half} and
\ref{thm:three-core-lower}. The last two rows are the earlier exact
large-cap theorem~\cite{Prior}, whose complete proofs are in that
unpublished companion and are not repeated here. The formulas agree at
shared endpoints.

Theorem~\ref{t:main} first determines the positive discrepancy for every
$a\ge1/2$. A competitive threshold forces two large coefficients and
an explicit residual variance obstruction. A four-coordinate argument
settles the corresponding positive side on $[1/\sqrt5,1/2]$.
The new lower-discrepancy argument expresses the law as a four-shift
mixture and bounds its strict central interval. Two finite rational-box
partitions and one three-box partition certify the remaining scalar
inequalities; their full analytic interface and reproducible enclosure
algorithm are given in Section~\ref{sec:lower-certificates}.
In fact the lower result is stronger than the full-distance formula
requires: $\mathfrak L(a)=c(\mathfrak G(a))$ for every
$a\ge1/\sqrt5$, and its optimizers are exactly the support-gap
optimizers, including all ties. Thus the deterministic classification
identifies every lower-discrepancy optimizer throughout that range.

The preceding package also established $\mathfrak K(a)$ for every
sufficiently small cap, with an unevaluated positive cutoff~\cite{Prior}.
That asymptotic proof remains in the companion. The deterministic gap
formula alone does not exclude competitive thresholds containing
positive central mass; the new probability reductions supply this
missing control on the two additional finite-cap intervals.

\subsection{Comparison with the literature}
Keller and Klein's work on Tomaszewski's conjecture includes sharp
normal-approximation questions for Rademacher sums~\cite{KK}. The
original small-coefficient research package~\cite{Original} supplies the
missing upper-concentration regime relative to Ivanisvili's October 2026
manuscript~\cite{Ivanisvili}. Claims using the completed comparison retain
the external verification boundary described in Section~\ref{sec:positive}.
Hollom and Portier determine the uncapped strict-tail profile used in our
new probability argument~\cite{HP}; our coefficient-cap optimization is
a different extremal problem. Their method already uses capped tail
bounds and conditioning on leading signs; those devices are not claimed
as new here.

Sharp coefficient stability has substantial precedents. De, Diakonikolas,
and Servedio bound the excess of $\E|X_b|$ above $1/\sqrt2$ by distance
from its two-flat minimizers~\cite{DDS}. Fang and Wang give a sharp
nonlinear profile near that set~\cite{FW}; Xie determines the optimal
global linear constant $(2-\sqrt2)^{3/2}/8$, attained at four-flat
coefficients~\cite{Xie}. These are moment inequalities. Our objective is
the minimum absolute signed sum, with odd cores and coefficient caps;
these displayed sharpness vectors have zero support gap.
Xie's Fourier--Haagerup argument differs from the finite-slice and
tail-repair geometry used here. Bell studies near-maximal Berry--Esseen
ratios for general independent arrays and their common-lattice and
two-point structure~\cite{Bell}; that objective and normalization also
differ from the capped central-gap problem.

The equal-core-plus-residual construction itself is not novel; it occurs
in earlier Rademacher extremal investigations, including Pinelis
\cite{Pinelis}. Nor do we claim independent novelty for the elementary
hypersimplex or $\ell_1/\ell_2$ inequalities used below. The contribution
is their sharp consequence for actual signed-sum gaps, across every cap,
with a complete near-optimizer and equality analysis, together with the
exact probability cap profile on the two additional intervals.
The primary-source comparison found no identical prior theorem, but it
is not a guarantee of priority.

\section{Majority structure and the exact support-gap optimum}\label{gap:section}

The extremal value and classification in this section are from the earlier
research package~\cite{Prior}; their short proof is included because it
supplies the geometry used in all the new stability theorems.


For a finite positive coefficient vector $b=(b_1,\ldots,b_N)$ define
\[
 g(b)=
 \min_{\varepsilon\in\{-1,1\}^N}
 \left|\sum_{i=1}^N b_i\varepsilon_i\right|,
\]
and, for $0<a\le1$, put
\[
 \mathfrak G(a)=
 \sup_{\substack{N<\infty,\ \sum_i b_i^2=1\\0<b_i\le a}}
 g(b).
\]
Zero coefficients may be added or removed freely; signs of coefficients
may be absorbed into the signs in the minimum.

\begin{theorem}[Exact support-gap optimum for every cap]\label{gap:main}
Let $n\ge1$ be odd and let
\[
 \frac1{\sqrt{n+2}}\le a\le\frac1{\sqrt n}.
\]
Then
\begin{equation}\label{gap:formula}
 \mathfrak G(a)=\max\left\{
 \frac1{\sqrt{n+2}},\ a-\sqrt{1-na^2}\right\}.
\end{equation}
The supremum is attained. Up to permutations, changes of signs, and
zero coefficients, all maximizing vectors are as follows.
If the first candidate is strictly larger, the only maximizer consists
of $n+2$ equal coefficients $1/\sqrt{n+2}$.
If the second candidate is strictly larger, the only maximizer consists
of $n$ coefficients equal to $a$ and one coefficient
$\sqrt{1-na^2}$, with a zero residual omitted.
At a tie, precisely both types maximize.

There is exactly one tie in the open interval
$(1/\sqrt{n+2},1/\sqrt n)$. Writing $\lambda=1/\sqrt{n+2}$,
its location is
\begin{equation}\label{gap:switch}
 a_n=\frac{\lambda+\sqrt{n+1-n\lambda^2}}{n+1}
     =\frac{1+\sqrt{n^2+2n+2}}{(n+1)\sqrt{n+2}}.
\end{equation}
Thus the first type maximizes below $a_n$ and the second above $a_n$.
The formulas agree at shared endpoints of adjacent odd-$n$ intervals,
so they specify the optimum and every maximizing vector for every
$0<a\le1$.
\end{theorem}

\subsection{A high-gap sum has a majority core}

The following structural fact does not require variance normalization.
For positive coefficients, a Hamming layer consists of the sign vectors
having a fixed number of positive signs.

\begin{lemma}[Majority core]\label{gap:majority}
Let $0<b_i\le a$ be finitely many positive coefficients. Suppose
$x>a/2$ and
\begin{equation}\label{gap:empty}
 \left|\sum_i b_i\varepsilon_i\right|\ge x
 \quad\text{for every sign vector }\varepsilon.
\end{equation}
Put
\[
 C=\{i:b_i\ge x\},\qquad k=|C|,\qquad
 \ell=\sum_{i\notin C}b_i.
\]
Then $k$ is positive and odd. The sign of the complete sum is the
majority sign of the $k$ coordinates in $C$ and is independent of
every coordinate outside $C$.
Moreover, for every core sign pattern $\eta\in\{-1,1\}^C$,
\begin{equation}\label{gap:coremargin}
 \left|\sum_{i\in C}b_i\eta_i\right|\ge x+\ell.
\end{equation}
In particular,
\begin{equation}\label{gap:average}
 x+\ell\le\frac1k\sum_{i\in C}b_i.
\end{equation}
\end{lemma}
\begin{proof}
Flipping a coordinate outside $C$ changes the full sum by
$2b_i<2x$. Such an edge cannot join a sum at least $x$ to one at
most $-x$. Hence its endpoints have the same sign. The cube of
outside coordinates is connected, so the sign is constant on every
fiber with the core coordinates fixed.

For a fixed core outcome $S=\sum_{i\in C}b_i\eta_i$, the two
residual endpoints are $S-\ell$ and $S+\ell$. They have the same
sign, and each has magnitude at least $x$. Thus either
$S-\ell\ge x$ or $S+\ell\le-x$, proving
\eqref{gap:coremargin}. In particular $k=0$ is impossible.

Swap a positive and a negative sign at two core coordinates. The
core sum changes by at most
$2|b_i-b_j|\le2(a-x)<2x$.
By \eqref{gap:coremargin}, this change cannot reverse its sign.
Any two sign patterns in the same Hamming layer are connected by
such swaps. Therefore the sign is constant on each layer.
As the number of positive signs increases, these layer signs are
nondecreasing: one may pass from any layer to the next by changing
a negative sign to a positive one, which strictly increases the sum.
Negating all signs reverses the sign and exchanges layers $j$ and
$k-j$.

If $k$ were even, its middle layer would have a constant sign equal
to its own negative, which is impossible. Thus $k$ is odd, and
monotonicity together with sign reversal forces the majority rule:
layers with at least $(k+1)/2$ positive signs are positive, and the
others are negative.

Average the core sum over the layer with exactly $(k+1)/2$ positive
signs. Each coordinate has average sign $1/k$, since it is positive
in the fraction $(k+1)/(2k)$ of the layer. The average sum is
therefore $k^{-1}\sum_{i\in C}b_i$.
Every sum on this layer is at least $x+\ell$ by
\eqref{gap:coremargin}, proving \eqref{gap:average}.
\end{proof}

\paragraph{A pair-deletion proof used in the formalization.}
The core margin also gives the stronger coordinatewise conclusion
$x+\ell\le b_i$ for every $i\in C$.
Indeed, let $y=x+\ell$ and choose two core coefficients $u,v$.
For every signing of the remaining coordinates, the two opposite
assignments to this pair give $S+(u-v)$ and $S-(u-v)$.
Both have absolute value at least $y$, whereas
$|u-v|\le a-x<x\le y$.
They therefore have the same sign and force
$|S|\ge y+|u-v|\ge y$.
Deleting any pair preserves the gap. Repeated deletion would give
$0\ge y>0$ if the core size were even; if it is odd, one may retain
any designated coefficient and conclude that it is at least $y$.
This proves oddness and the coordinatewise bound directly, without
an assumed majority rule. Summing their squares gives
$k(x+\ell)^2\le1-V$ in the normalized case.

\begin{corollary}[Variance form]\label{gap:variance}
Under the hypotheses of Lemma~\ref{gap:majority}, assume in addition
$\sum_i b_i^2=1$ and put $V=\sum_{i\notin C}b_i^2$.
Then
\begin{equation}\label{gap:variancebound}
 x\le\sqrt{\frac{1-V}{k}}-\sqrt V.
\end{equation}
If $ka^2\le1$, this further implies
\begin{equation}\label{gap:capbound}
 x\le a-\sqrt{1-ka^2}.
\end{equation}
\end{corollary}
\begin{proof}
One has $\ell\ge\sqrt V$, and Cauchy--Schwarz gives
$k^{-1}\sum_{i\in C}b_i\le\sqrt{(1-V)/k}$.
Combine these inequalities with \eqref{gap:average}.
Also $V\ge1-ka^2$, and the function
$v\mapsto\sqrt{(1-v)/k}-\sqrt v$ decreases strictly on $[0,1]$.
Substituting its smallest permitted argument gives
\eqref{gap:capbound}.
\end{proof}

\subsection{Exact optimization and equality}

\begin{proof}[Proof of Theorem~\ref{gap:main}]
Put $\lambda=1/\sqrt{n+2}$ and $g=a-\sqrt{1-na^2}$.
The $(n+2)$-flat vector is admissible and has gap $\lambda$.
Whenever $g\ge\lambda$, the residual $r=\sqrt{1-na^2}$ satisfies
$0\le r\le a-\lambda<a$; hence the saturated vector of $n$
coefficients $a$ and the residual $r$ is admissible. Since $n$ is
odd and $r<a$, its minimum absolute signed sum is $a-r=g$.
Thus the right side of \eqref{gap:formula} is attained.

To prove the upper bound and classify equality, take an arbitrary
admissible vector and write $x=g(b)$.
If $x<\lambda$ there is nothing to prove. Otherwise
\[
 x\ge\lambda>\frac a2,
\]
because $a\le1/\sqrt n$ and $4n>n+2$ for $n\ge1$.
Apply Lemma~\ref{gap:majority} with the threshold $x$.
The number of core coefficients satisfies
$kx^2\le\sum_i b_i^2=1$, so $k\le n+2$.
If $k=n+2$, equality is forced throughout
\[
 1\ge\sum_{i\in C}b_i^2\ge(n+2)x^2\ge(n+2)\lambda^2=1.
\]
Thus $x=\lambda$, all core coefficients equal $\lambda$, and
there is no nonzero residual. This is exactly the $(n+2)$-flat type.

In the remaining case $k<n+2$. Since $k$ is odd, $k\le n$.
In particular $ka^2\le na^2\le1$, so
Corollary~\ref{gap:variance} gives
\begin{equation}\label{gap:chain}
 x\le a-\sqrt{1-ka^2}\le a-\sqrt{1-na^2}=g.
\end{equation}
This proves \eqref{gap:formula}.

If a maximizer is not the $(n+2)$-flat type, then it belongs to the
second case and $x=\max\{\lambda,g\}\le g$. Hence $x=g$ and
all inequalities in \eqref{gap:chain} and its proof are equalities.
The last inequality is strict for $k<n$, so $k=n$.
Strict decrease in the variance substitution gives
$V=1-na^2$. Since all $n$ core coefficients are at most $a$,
their squared sum $na^2$ forces all of them to equal $a$.
Equality $\ell=\sqrt V$ forces at most one nonzero residual
coefficient: for nonnegative numbers,
$(\sum b_i)^2=\sum b_i^2$ holds only when at most one is nonzero.
That residual is $\sqrt{1-na^2}$, or is absent when this value
is zero. Conversely this vector has already been shown to attain
the active second candidate. The asserted classification follows,
including at a tie.

Finally $a\mapsto a-\sqrt{1-na^2}$ increases strictly on the
interval in the theorem. At its left endpoint it equals
$(1-\sqrt2)\lambda<\lambda$, and at its right endpoint it is
$1/\sqrt n>\lambda$.
Solving $a-\sqrt{1-na^2}=\lambda$ gives
\eqref{gap:switch}; the positive solution is the unique one in
the interval. At a shared endpoint $a=1/\sqrt n$, the result
is the $n$-flat vector, matching the neighboring interval's
constant branch. This also covers $n=1$ and $a=1$.
\end{proof}


\section{A sharp global inverse theorem for large support gaps}
\label{sec:sharp-gap-rigidity}

For a nonzero finite real vector $b$, write
\[
 g(b)=\min_{\varepsilon_i\in\{-1,1\}}\left|\sum_i b_i\varepsilon_i\right|,
 \qquad m(b)=\max_i|b_i|.
\]
Signs of coefficients may be absorbed into the independent signs, so the
proofs below use nonnegative coefficients and allow zeros.  For a nonempty
set $C$ of cardinality $k$, let $e_C=k^{-1/2}\mathbf1_C$.

\begin{theorem}[Global sharp gap inverse]\label{thm:sharp-gap-inverse}
Suppose $\sum_i b_i^2=1$, $|b_i|\le a\le1$, and
$a/2<x\le g(b)$. The cap need not be attained and $x$ may be any
certified lower bound for the actual gap. Let
\[
 C=\{i:|b_i|\ge x\},\quad k=|C|,\quad \delta=a-x,
 \quad \rho=\bigl\|(|b_i|)_i-e_C\bigr\|_2.
\]
Then $k$ is positive and odd, the sign of the signed sum is the majority
sign on $C$, and
\begin{equation}\label{eq:sharp-gap-inverse}
 \rho^2\le 2-2\sqrt{1-\delta^2}.
\end{equation}
Equivalently,
$\rho\sqrt{1-\rho^2/4}\le\delta$.
Equality holds if and only if, after signs, permutations, and zeros are
discarded, the vector consists of $k$ copies of $a$ and one residual
coefficient $\delta$; the residual is omitted when $\delta=0$.
Necessarily $ka^2+\delta^2=1$ and $\delta<a/2$; equality also forces
$a=m(b)$ and $x=g(b)$.

For every odd $k\ge1$ and every $0\le r<1/\sqrt{4k+1}$, the vector
with $k$ coefficients $\sqrt{(1-r^2)/k}$ and one coefficient $r$
satisfies the hypotheses and attains equality.  In particular the leading
linear coefficient in $\rho\le(1+o(1))\delta$ cannot be reduced.
\end{theorem}

The inequality is valid throughout the high-gap regime, with no asymptotic
smallness assumption and no numerical integration.  Its equality family
is specified explicitly; the assertion is not that this right side is an
attained optimum at every arbitrarily prescribed value of $\delta$.

\subsection{A sharp finite-slice inequality}

\begin{lemma}\label{lem:sharp-slice-norm}
Let $k=2h+1\ge3$, let $d_1,\ldots,d_k$ be real numbers with
$\sum_i d_i=0$, and put
\[
 s=\left(\sum_i d_i^2\right)^{1/2},\quad
 M=\max_i d_i,\quad
 A=\max_{|J|=h}\sum_{i\in J}d_i,\quad
 \gamma=\sqrt{\frac{k-1}{k}}.
\]
Then
\begin{equation}\label{eq:two-slice-bounds}
 2A\ge\gamma s,\qquad M\ge\frac{s}{k\gamma},
 \qquad s\le\gamma(M+2A).
\end{equation}
For $s>0$, equality in either of the first two inequalities occurs
precisely when $k-1$ entries are equal and positive, and the last entry
is minus their sum.  The same statement holds for equality in the third.
\end{lemma}
\begin{proof}
First, whenever $0\le z_i\le1$ and $\sum_i z_i=h$, one has
\begin{equation}\label{eq:fractional-slice}
 \sum_i d_i z_i\le A.
\end{equation}
Indeed, a nonintegral such vector has at least two fractional coordinates.
Choose two, say $i,j$, with $d_i\ge d_j$, and increase $z_i$ while
decreasing $z_j$ by the same amount until at least one becomes integral.
This preserves feasibility and does not decrease the objective. Repeating
terminates at the indicator of an $h$-subset and proves the claim.

Assume $s>0$. Since $d_i=-\sum_{j\ne i}d_j$, Cauchy--Schwarz gives
$d_i^2\le(k-1)(s^2-d_i^2)$, hence $|d_i|\le\gamma s$.
Set
\[
 r=\frac{h}{\sqrt{k(k-1)}}=\frac{h}{k\gamma},
 \qquad z_i=\frac hk+r\frac{d_i}{s}.
\]
Then $0\le z_i\le2h/k<1$ and $\sum_i z_i=h$.
Equation~\eqref{eq:fractional-slice} yields
$A\ge rs$, so $2A\ge\gamma s$.

If equality holds, some $z_j$ must be zero. Otherwise all coordinates
are strictly between zero and one, and a small transfer from a coordinate
with smaller $d_i$ to one with larger $d_i$ strictly increases the
objective, contradicting equality in~\eqref{eq:fractional-slice}.
Thus $d_j=-\gamma s$, and equality in the displayed Cauchy--Schwarz
inequality forces the other $k-1$ entries to be equal and positive.
Conversely that vector attains the bound.

For the second inequality put $u_i=M-d_i\ge0$. Then
\[
 s^2+kM^2=\sum_i u_i^2\le\left(\sum_i u_i\right)^2=k^2M^2.
\]
Thus $M\ge s/\sqrt{k(k-1)}$.  Equality for $s>0$ requires at most
one $u_i$ to be nonzero, which is the same vector shape.
Finally $\gamma+1/(k\gamma)=1/\gamma$, so adding the first two
inequalities proves the third and its equality statement. The case
$s=0$ is immediate.
\end{proof}

\subsection{From a genuine support gap to the coefficient distance}

\begin{proof}[Proof of Theorem~\ref{thm:sharp-gap-inverse}]
Replace $b_i$ by $|b_i|$. Put
\[
 \ell=\sum_{i\notin C}b_i,\qquad V=\sum_{i\notin C}b_i^2.
\]
Lemma~\ref{gap:majority} and its pair-deletion argument give positive
odd $k$, the majority rule, and $b_i\ge x+\ell$ for every core
coordinate. In particular the least positive core sum lies on the
nearest majority layer.

Let $\beta=k^{-1}\sum_{i\in C}b_i$ and
$T=\sum_{i\in C}(b_i-\beta)^2$. The actual maximum lies in $C$.
The least positive core sum occurs on a nearest majority layer, so, with
$h=(k-1)/2$ and $A=\max_{|J|=h}\sum_{i\in J}(b_i-\beta)$,
\[
 x+\ell\le\beta-2A,
 \qquad (a-\beta)+2A\le\delta-\ell.
\]
In particular $0\le\ell\le\delta$.
If $k\ge3$, Lemma~\ref{lem:sharp-slice-norm}, with
$\max_{i\in C}(b_i-\beta)\le a-\beta$, gives
\begin{equation}\label{eq:gap-variance-defect}
 T\le\frac{k-1}{k}(\delta-\ell)^2.
\end{equation}
For $k=1$, $T=0$, so the same formula remains true.
Since $V\le\ell^2$,
\begin{equation}\label{eq:total-defect}
 T+V\le\frac{k-1}{k}(\delta-\ell)^2+\ell^2\le\delta^2.
\end{equation}
The last inequality follows for $0\le\ell\le\delta$ by expansion;
if $\delta>0$ it is strict unless $\ell=\delta$.

Normalization gives $k\beta^2+T+V=1$. Since $\beta>0$,
\[
 \rho^2=2-2\sqrt{k}\,\beta
       =2-2\sqrt{1-(T+V)}
       \le2-2\sqrt{1-\delta^2}.
\]
Here $0\le\delta<a/2\le1/2$, so every square root is well defined.
The identity $T+V=\rho^2-\rho^4/4$ also gives the equivalent form.

If $\delta>0$, equality forces $\ell=\delta$, $T=0$, and
$V=\ell^2$. Hence the core is flat. Its value $\beta$ satisfies
$a=x+\delta=x+\ell\le\beta\le a$, so the core value is $a$,
and precisely one nonzero tail coefficient equals $\delta$.
If $\delta=0$, the same inequalities give $T=V=0$ and $\beta=a$, so
the vector is odd flat. Conversely, for an odd flat core of coefficient
$a$ and one residual $r<a/2$, the actual gap is $a-r$, the core selected by its gap
is exactly the flat core, and
$\rho^2=2-2\sqrt{1-r^2}$. Choosing
$a=\sqrt{(1-r^2)/k}$ proves the asserted equality families and sharpness.
\end{proof}

\paragraph{Sharpness of the structural threshold.}
The strict threshold $g>m/2$ cannot be relaxed to $g\ge m/2$.
For the normalized vector $(2,1,1,1)/\sqrt7$, the gap is $1/\sqrt7$
and the maximum is $2/\sqrt7$, while its gap-defined core has even
cardinality four. The large-gap core theorem therefore has the exact
threshold $1/2$.

\subsection{A complementary sharp bound relative to a fixed flat dimension}

\begin{proposition}\label{prop:fixed-flat-gap}
In the actual-gap case $x=g(b)=g$ of
Theorem~\ref{thm:sharp-gap-inverse}, if $k\ge3$,
\[
 g\le\frac{1-\rho^2/2}{\sqrt{k}}
       -\sqrt{\frac{k-1}{k}}\,\rho\sqrt{1-\rho^2/4}.
\]
Equivalently,
\[
 \rho^2\le
 2-\frac2{\sqrt{k}}\left(g+\sqrt{(k-1)(1-g^2)}\right).
\]
Equality holds precisely for $k-1$ equal high coefficients and one low
coefficient, with no tail. These vectors are
\[
 \left(\underbrace{\sqrt{\frac{1-g^2}{k-1}},\ldots,
 \sqrt{\frac{1-g^2}{k-1}}}_{k-1},g\right),
 \qquad \frac1{\sqrt{4k-3}}<g\le\frac1{\sqrt{k}}.
\]
For a prescribed cap $a$, the original conditions
$\sqrt{(1-g^2)/(k-1)}\le a$ and $g>a/2$ must also hold.
At $a=1/\sqrt{k}$ the nonflat equality family is inadmissible;
on an interior flat phase it is admissible for all sufficiently small
positive deficits.
\end{proposition}
\begin{proof}
The first slice inequality gives
$g\le\beta-\gamma\sqrt T-\ell$.
Since $\ell\ge\sqrt V$ and $0<\gamma<1$,
$\gamma\sqrt T+\sqrt V\ge\gamma\sqrt{T+V}$, with equality only
if $V=0$. Substitute
$\beta=(1-\rho^2/2)/\sqrt{k}$ and
$T+V=\rho^2-\rho^4/4$.
For the equivalent form put
$\cos\theta=\sqrt{k}\beta=1-\rho^2/2$ and
$\cos\psi=1/\sqrt{k}$, with $\theta,\psi\in[0,\pi/2)$.
The first inequality reads $g\le\cos(\theta+\psi)$; its positive
right side implies $\theta+\psi<\pi/2$, so it can be inverted.
This yields
$\cos\theta\ge[g+\sqrt{(k-1)(1-g^2)}]/\sqrt{k}$.
Equality follows from the exact equality cases in
Lemma~\ref{lem:sharp-slice-norm}; the high-gap condition is precisely
$g>\frac12\sqrt{(1-g^2)/(k-1)}$.
\end{proof}


\section{Uniform sharp stability of the exact cap optimum}
\label{sec:all-cap-gap-stability}

Let $n\ge1$ be odd and $1/\sqrt{n+2}\le a\le1/\sqrt n$. Put
\[
 \lambda=(n+2)^{-1/2},\quad \mu=(n+4)^{-1/2},\quad
 r=\sqrt{1-na^2},\quad s=a-r,\quad G=\max\{\lambda,s\},
 \quad C_* =\sqrt{2+\sqrt2}.
\]
Let $v_f$ denote the $(n+2)$-flat vector, and, when $r\le a$,
let $v_s=(a^n,r)$, with zero entries omitted. All coefficient vectors
in distances are sorted in decreasing magnitude and padded with zeros.

\begin{theorem}[Sharp linear stability across all cap bands]
\label{thm:all-cap-gap-stability}
Let $b$ be any finite normalized coefficient vector with $|b_i|\le a$.
Write $x=g(b)$ and $\Delta=G-x$. If
\begin{equation}\label{eq:stability-trigger}
 x>\max\{a/2,\mu\},
\end{equation}
then the actual core $\{i:|b_i|\ge x\}$ has size $n$ or $n+2$.
More precisely:
\begin{enumerate}
\item If its size is $n+2$, then $x\le\lambda$ and
\[
 \|b-v_f\|_2\le
 \sqrt{\frac{n+2}{n+1}}(\lambda-x)\le C_*\Delta.
\]
\item If its size is $n$, then $r<a/2$, the saturated candidate is
admissible, $x\le s$, and
\[
 \|b-v_s\|_2\le C_*(s-x)\le C_*\Delta.
\]
\end{enumerate}
Consequently the distance to the admissible members of these two candidate types is at
most $C_*\Delta$. If the branches are unequal and
$\Delta<|\lambda-s|$, the candidate selected above is necessarily the
unique maximizing type. At a branch tie both candidates are maximizers.

The constant $C_*$ is optimal, even on a subinterval with a unique
maximizer and under the strict trigger~\eqref{eq:stability-trigger}.
The endpoint $\mu$ in the core-count trigger cannot be included in a
uniform linear estimate of this form.
\end{theorem}

Unlike a bound that chooses a unique branch continuously in $a$, this
statement remains valid through the switching caps. The need to keep
both competing types near a switch is real: on either side of a switch,
the losing candidate has deficit tending to zero while its distance to
the unique winner tends to the nonzero distance between the two shapes.

\subsection{Sharp concentration of a nonnegative tail}

\begin{lemma}\label{lem:sharp-tail-concentration}
Let $t$ be a finite nonnegative vector, $\sigma=\|t\|_2$,
$\ell=\|t\|_1$, and let $j$ be a largest coordinate, after adding a zero
coordinate if the index set is empty. Then
\[
 \|t-\sigma e_j\|_2\le C_*(\ell-\sigma).
\]
If $\ell>\sigma$, equality holds if and only if exactly two coordinates
are nonzero and they are equal. If $\ell=\sigma$, both sides vanish.
\end{lemma}
\begin{proof}
Only $\sigma>0$ and $\ell>\sigma$ need consideration. Normalize
$\sigma=1$, and put $z=t_j\in(0,1)$. The squared distance is $2(1-z)$.
If $z\le1/\sqrt2$, then $1=\sum_i t_i^2\le z\ell$, and hence
\[
 \frac{2(1-z)}{(\ell-1)^2}\le\frac{2z^2}{1-z}\le2+\sqrt2.
\]
The last expression increases strictly with $z\in(0,1)$, and its value
at $1/\sqrt2$ is $2+\sqrt2$.
If $z\ge1/\sqrt2$, then
$\ell\ge z+\sqrt{1-z^2}$. Put
$u=\sqrt{(1-z)/(1+z)}\in(0,\sqrt2-1]$. Direct substitution gives
\[
 \frac{\sqrt{2(1-z)}}{z+\sqrt{1-z^2}-1}
 =\frac{\sqrt{1+u^2}}{1-u}\le\sqrt{2+\sqrt2}.
\]
The middle expression increases strictly with $u$, as is seen by
differentiation. In either case equality forces $z=1/\sqrt2$.
The first case further requires every nonzero $t_i$ to equal $z$;
the second requires at most one further nonzero coordinate.
Both give exactly two equal nonzero coordinates. The converse is direct.
\end{proof}

\subsection{The saturated branch}


\begin{lemma}[Repair of an almost saturated core]\label{lem:saturated-repair}
Let $b$ be normalized and nonnegative, with coefficients at most $a$,
actual gap $x>a/2$, and core size $k$ satisfying $ka^2\le1$.
Write
\[
 r=\sqrt{1-ka^2},\quad e=a-r-x,
 \quad \ell=\|b_{C^c}\|_1,\quad \sigma=\|b_{C^c}\|_2,
\]
and put $q=\ell-\sigma$, $D=\sigma-r$, and $\kappa=a-x-\ell$.
Then $q,D,\kappa\ge0$ and $e=q+D+\kappa$.
If $j$ is a largest tail coordinate, with a zero coordinate added when
necessary, then
\begin{equation}\label{eq:saturated-repair}
 \operatorname{dist}(b,(a^k,r))
 \le \|b_{C^c}-\sigma e_j\|_2+\kappa+D.
\end{equation}
\end{lemma}
\begin{proof}
The cap and normalization give $\sigma\ge r$, and the core lemma gives
$\ell\le a-x$ and $x+\ell\le b_i\le a$ on the core. These prove
the asserted signs and identity. Put $d_i=a-b_i$ on the core,
$S=\sum d_i$, and $Q=\sum d_i^2$. Then $0\le d_i\le\kappa$,
so $Q\le\kappa S$ and
\[
 \sigma^2-r^2=2aS-Q\ge(2a-\kappa)S.
\]
If $\sigma+r=0$, normalization forces every core coefficient to equal
$a$, and the assertion is immediate. Otherwise
\[
 D=\frac{\sigma^2-r^2}{\sigma+r}\ge S,
 \qquad
 \sqrt{Q+D^2}\le\sqrt{\kappa D+D^2}\le\kappa+D,
\]
because $\sigma+r\le2\ell=2(a-\kappa-x)<2a-\kappa$.
First merge the tail to a coefficient $\sigma$ at $j$. Changing the core
to $k$ copies of $a$ and this residual from $\sigma$ to $r$ then costs
$\sqrt{Q+D^2}$ in Euclidean distance. The triangle inequality proves
\eqref{eq:saturated-repair}. Since core coefficients exceed every tail
coefficient, this alignment gives the distance under our sorting convention.
\end{proof}

\begin{lemma}\label{lem:saturated-gap-stability}
Under the preceding hypotheses,
\[
 \operatorname{dist}(b,(a^k,r))\le C_*e.
\]
For $e>0$, equality holds precisely when the core consists of $k$ copies
of $a$ and the tail consists of two equal coefficients $r/\sqrt2$;
necessarily $e=(\sqrt2-1)r$.
\end{lemma}
\begin{proof}
The tail concentration lemma and~\eqref{eq:saturated-repair} give
\[
 \operatorname{dist}(b,(a^k,r))\le C_*q+(e-q)\le C_*e.
\]
Since $C_*>1$, equality with $e>0$ forces $q=e$, hence $D=\kappa=0$.
The core is therefore saturated and $\sigma=r$. Equality in the tail
lemma gives exactly two equal residual coefficients. Their sum is
$\sqrt2r$, yielding the stated deficit. Conversely this vector attains
the bound whenever its actual gap $a-\sqrt2r$ exceeds $a/2$.
When $e=0$, all the nonnegative quantities vanish and the vector is the
saturated comparator itself.
\end{proof}

\subsection{The flat branch and the exact core-count trigger}

\begin{proof}[Proof of Theorem~\ref{thm:all-cap-gap-stability}]
Absorb signs and discard zeros. Let $k$ be the positive odd core size,
and let $\ell,V$ be the tail sum and variance. The core lemma gives
$\ell\le a-x<a/2$. Hence
\[
 1=\sum_C b_i^2+V\le ka^2+\ell^2<(k+1/4)a^2\le(k+1/4)/n.
\]
Thus $k\ge n$. Also $kx^2\le1$ and $x>\mu$ imply $k<n+4$.
Oddness gives exactly $k=n$ or $k=n+2$.

For $k=n$, use Lemma~\ref{lem:saturated-gap-stability}. It also gives
$r\le\sigma\le\ell<a/2$, so the saturated candidate is admissible.

For $k=n+2$, put $\gamma=\sqrt{(k-1)/k}$ and $\lambda=1/\sqrt{k}$.
Proposition~\ref{prop:fixed-flat-gap} yields, with
$\rho=\|b-v_f\|_2$ and $S=\rho^2-\rho^4/4$,
\[
 x\le\lambda\sqrt{1-S}-\gamma\sqrt S.
\]
Since $x>0$, this implies $S<\lambda^2$. Moreover
$\rho^2=2-2\sqrt{1-S}\le2S$, so $\gamma\rho/2<\lambda$.
Using $\sqrt{1-\rho^2/4}\ge1-\rho^2/4$ gives
\[
 \lambda-x\ge\lambda\rho^2/2+\gamma\rho\sqrt{1-\rho^2/4}
 \ge\gamma\rho+\frac{\rho^2}{2}(\lambda-\gamma\rho/2)
 \ge\gamma\rho.
\]
This proves the stated flat estimate. Its leading coefficient $1/\gamma$ is sharp on the interior of a
flat phase by the admissible one-low-coordinate equality family in
Proposition~\ref{prop:fixed-flat-gap}.

Every core-$n$ vector has gap at most $s$, and every core-$(n+2)$ vector
has gap at most $\lambda$. Therefore a deficit below $|\lambda-s|$
excludes the losing core size. The exact cap formula identifies the
remaining candidate as the actual optimizer.

To prove sharpness of $C_*$, take $n=1$, let $r>0$ tend to zero, and set
$a=\sqrt{1-r^2}$. Then the unique optimizer is $(a,r)$.
For $b=(a,r/\sqrt2,r/\sqrt2)$,
\[
 x=a-\sqrt2r,\quad \Delta=(\sqrt2-1)r,
 \quad\operatorname{dist}(b,(a,r))=r\sqrt{2-\sqrt2}=C_*\Delta.
\]
For sufficiently small $r$, the strict trigger holds, the losing
three-flat vector stays a positive distance away, and
$\Delta<|\lambda-s|$. Thus the constant is sharp even for distance to
the unique global optimizer.

Finally take $a=1/\sqrt{n+2}$ and let $b$ be $(n+4)$-flat. Its gap is
exactly $\mu$, the saturated candidate is inadmissible, and $G=\lambda$.
Writing $k=n+2\ge3$, the squared ratio of its distance from the $k$-flat
optimizer to its deficit is
\[
 \frac{2-2\sqrt{k/(k+2)}}{(k^{-1/2}-(k+2)^{-1/2})^2}
 =k\bigl(k+2+\sqrt{k(k+2)}\bigr)>C_*^2.
\]
The ratio is unbounded as $k\to\infty$. This proves the endpoint
obstruction and completes the proof.
\end{proof}

\paragraph{Interpretation and limit.}
For $n\ge3$, $\mu>a/2$ throughout the band, and the trigger is simply
$\Delta<G-\mu$. On the flat branch this allowable deficit is
$\lambda-\mu\asymp a^3$. The theorem is a sharp quantitative extension
of the all-cap support-gap classification, valid without an unevaluated
small-cap cutoff. It does not resolve the remaining probability
optimization at thresholds containing positive central mass.

\section{The exact local stability modulus on the saturated branch}
\label{sec:exact-saturated-stability}

The optimal linear constant is the endpoint of an exact nonlinear
stability modulus. For $r>0$ and $0\le e\le(\sqrt2-1)r$, define
\[
 F_r(e)=\sqrt{r\left(r-e-\sqrt{r^2-2re-e^2}\right)}.
\]
Both square roots are nonnegative on this interval. Indeed
$(r-e)^2-(r^2-2re-e^2)=2e^2$ and $e<r$.

\begin{theorem}[Exact saturated stability modulus]
\label{thm:exact-saturated-stability}
Let $b$ be a finite normalized nonnegative vector with coefficients at
most $a\le1$, actual gap $x>a/2$, and actual core of odd size $k$.
Suppose $ka^2<1$, put
\[
 r=\sqrt{1-ka^2}>0,\qquad e=a-r-x,
\]
and assume $e\le(\sqrt2-1)r$. Then $e\ge0$ and
\begin{equation}\label{eq:exact-saturated-stability}
 \operatorname{dist}(b,(a^k,r))\le F_r(e).
\end{equation}
Equality holds precisely when the core consists of $k$ copies of $a$
and the tail has the two coefficients
\begin{equation}\label{eq:two-residual-modulus}
 t_\pm=\frac{r+e\pm\sqrt{r^2-2re-e^2}}2,
\end{equation}
with the zero coefficient omitted at $e=0$.
Conversely these coefficients attain the bound whenever the high-gap
condition $a-r-e>a/2$ holds.

In particular $F_r(e)/e$ increases strictly from $1$ to
$\sqrt{2+\sqrt2}$ as $e/r$ increases from $0$ to $\sqrt2-1$.
\end{theorem}

When the saturated type is the unique cap optimizer and the deficit is
below the branch separation, this gives the exact coefficient-distance
modulus for near optimizers in the stated deficit range. The two small
residuals in~\eqref{eq:two-residual-modulus} are forced at equality;
the second residual cannot be replaced by an arbitrary small diffuse tail.

\begin{lemma}[Exact near-axis tail modulus]\label{lem:exact-tail-modulus}
Let $t$ be a finite nonnegative vector with $\sigma=\|t\|_2>0$,
$\ell=\|t\|_1$, and $q=\ell-\sigma\le(\sqrt2-1)\sigma$.
If $j$ is a largest coordinate, then
\[
 \|t-\sigma e_j\|_2\le F_\sigma(q).
\]
For $q>0$, equality holds precisely when at most two coordinates are
nonzero. Their values are obtained by replacing $(r,e)$ with $(\sigma,q)$
in~\eqref{eq:two-residual-modulus}. For $q=0$ the vector is one-sparse.
\end{lemma}
\begin{proof}
Put $M=t_j$. Since $\sigma^2\le M\ell$ and $\ell\le\sqrt2\sigma$,
one has $M\ge\ell/2$. Also
\[
 \sigma^2-M^2\le(\ell-M)^2,
\]
with equality exactly when at most one coordinate besides $j$ is
nonzero. Solving this quadratic inequality on $M\ge\ell/2$ gives
\[
 M\ge\frac{\ell+\sqrt{2\sigma^2-\ell^2}}2.
\]
The squared distance is $2\sigma(\sigma-M)$. Substituting the last
bound yields $F_\sigma(q)^2$, with exactly the stated equality cases.
\end{proof}

\begin{lemma}[Monotonicity of the modulus]\label{lem:modulus-monotonicity}
The function $R(\alpha)=F_1(\alpha)/\alpha$ is strictly increasing on
$0<\alpha\le\sqrt2-1$, has limit $1$ at zero, and has endpoint value
$\sqrt{2+\sqrt2}$. Consequently, whenever $0\le q\le e\le(\sqrt2-1)r$
and $\sigma\ge r>0$,
\[
 F_\sigma(q)\le F_r(q),\qquad F_r(q)+(e-q)\le F_r(e).
\]
If $e>0$, the second inequality is strict unless $q=e$.
\end{lemma}
\begin{proof}
For $0<\alpha\le\sqrt2-1$, rationalization gives
\[
 R(\alpha)^2=
 \frac{2}{1-\alpha+\sqrt{1-2\alpha-\alpha^2}}.
\]
The denominator decreases strictly from $2$ to $2-\sqrt2$, proving
the monotonicity and the stated endpoint values.
Since $F_\sigma(q)=qR(q/\sigma)$, the first consequence follows.
For the second, if $q>0$,
\[
 F_r(e)-F_r(q)
 =(e-q)R(e/r)+q\bigl(R(e/r)-R(q/r)\bigr)\ge e-q.
\]
The same conclusion for $q=0$ follows from $R(e/r)\ge1$.
If $e>q$ and $e>0$, then $R(e/r)>1$, so the inequality is strict.
\end{proof}

\begin{proof}[Proof of Theorem~\ref{thm:exact-saturated-stability}]

Use the nonnegative quantities $q,D,\kappa$ and the identity
$e=q+D+\kappa$ from Lemma~\ref{lem:saturated-repair}. Since $r>0$,
the tail norm $\sigma$ is positive. The exact tail lemma and the modulus
lemma yield
\[
 \operatorname{dist}(b,(a^k,r))
 \le F_\sigma(q)+(e-q)
 \le F_r(q)+(e-q)\le F_r(e).
\]
Their hypotheses hold because
$q\le e\le(\sqrt2-1)r\le(\sqrt2-1)\sigma$.

If $e>0$, equality forces $q=e$ by strictness in the modulus lemma.
The identity $e=q+D+\kappa$ then gives $D=\kappa=0$.
Every core coefficient equals $a$, the tail variance is $r^2$, and
equality in the exact tail lemma gives~\eqref{eq:two-residual-modulus}.
If $e=0$, the same nonnegative decomposition gives $q=D=\kappa=0$,
so the vector is the saturated optimizer itself.
Conversely, the displayed two residuals have sum $r+e$ and squared sum
$r^2$. As $k$ is odd and $r+e<a/2$, the actual gap is $a-r-e$;
both residuals lie below this gap. The exact tail calculation gives
equality in~\eqref{eq:exact-saturated-stability}.
\end{proof}

\paragraph{Local expansion.}
For fixed $r>0$ and $e/r\to0$,
\[
 F_r(e)=e+\frac{e^2}{2r}+O(e^3/r^2).
\]
Thus the sharp leading linear coefficient on every fixed saturated
branch is $1$. The larger universal coefficient
$\sqrt{2+\sqrt2}$ is needed when the entire residual splits into two
equal pieces, at $e=(\sqrt2-1)r$.

\section{Exact positive discrepancy for every cap at least one half}\label{sec:positive}

Let $X=\sum_i b_i\varepsilon_i$, $\sum_i b_i^2=1$, and $|b_i|\le a$,
with finitely many independent uniform signs. Write $c(x)=\Phi(x)-1/2$,
$U_X(x)=\frac12\Pp(|X|\le x)-c(x)$ for $x\ge0$, and
\[
 \mathfrak D(a)=\sup_{N,b}\sup_{x\ge0}U_X(x).
\]

This section uses two stated analytic inputs. First, the completed
Keller--Klein comparison for normalized Rademacher sums is
\begin{equation}\label{eq:KK-input}
 d_K(Z,N(0,1))\le c\!\left(\max_i|b_i|\right).
\end{equation}
Its small-coefficient part is proved in the original package
\cite{Original}; the passage to every coefficient size uses the external
computer-assisted theorem of Ivanisvili~\cite{Ivanisvili}. The expanded
companion~\cite{Prior} gives another reconstruction with its own
explicitly recorded external inputs. We use~\eqref{eq:KK-input} as a
theorem input here and do not claim to reprove its external certificates.

Second, the published theorem of Hollom and Portier~\cite{HP} gives
the following strict-tail lower bounds:
\begin{equation}\label{eq:HP-input}
 \Pp(|Z|>t)\ge f(t),\qquad
\begin{array}{c|ccccc}
t &[0,1/\sqrt7)&[1/\sqrt7,1/\sqrt5)&[1/\sqrt5,1/\sqrt3)&
[1/\sqrt3,\sqrt{2/3})&[\sqrt{2/3},1)\\\hline
f(t)&1/2&29/64&3/8&1/4&7/32.
\end{array}
\end{equation}
Their theorem is stated for non-strict tails. For a fixed $t$ in a listed
half-open interval, apply it at $t+\eta$ inside that interval and let
$\eta\downarrow0$. Continuity of probability from below gives exactly
\eqref{eq:HP-input}, including the left endpoints. The external
computer-assisted part of their proof is not replayed in this package.

The probability theorem below is thus conditional on these precise
external analytic inputs in the verification sense. The deterministic
theorems in the previous sections have no such dependencies.

\begin{theorem}\label{t:main}
For $1/2\le a\le1/\sqrt2$, put $r=\sqrt{1-2a^2}$. Then
\[
 \mathfrak D(a)=\max\{3/16,\ 1/4-c(r)\}.
\]
Let $q>0$ be the unique solution of $c(q)=1/16$, and put
$a_0=\sqrt{(1-q^2)/2}$. For $1/2\le a<a_0$, the only optimizer
is the four-flat vector, at threshold zero. For $a_0<a\le1/\sqrt2$,
the only optimizer is $(a,a,r)$, at threshold $r$, omitting a zero
residual. At $a=a_0$, precisely both types optimize. These statements
are up to permutations, sign changes, and zero padding.
For $a\ge1/\sqrt2$, $\mathfrak D(a)=1/4$, with only the equal pair,
at threshold zero.
\end{theorem}


\subsection{Elementary localization estimates}
We use the scalar inequalities
\begin{equation}\label{eq:scalar-inputs}
 \frac{199}{500}<\varphi(0)<\frac25,\quad
 c(1/\sqrt7)>\frac9{64},\quad c(1/\sqrt3)>\frac{43}{200},
 \quad 1-\Phi(1)<\frac7{40},\quad q<\frac4{25}.
\end{equation}
Here are elementary proofs. The classical bounds $223/71<\pi<22/7$
imply the first pair. Since $e^{-u}\ge1-u$ for $u\ge0$,
\[
 c(t)>\frac{199}{500}\left(t-\frac{t^3}{6}\right)\quad(0<t\le1).
\]
For $t=1/\sqrt7$ use $t>37/100$ and $t^2=1/7$; for
$t=1/\sqrt3$ use $t>577/1000$ and $t^2=1/3$. The resulting rational
lower bounds exceed $9/64$ and $43/200$. At $t=1$, the lower bound
exceeds $13/40$, proving the tail assertion. Finally
\[
 c(4/25)>\frac{199}{500}
 \left(\frac4{25}-\frac{(4/25)^3}{6}\right)>\frac1{16},
\]
so strict monotonicity of $c$ gives $q<4/25$.
All remaining numerical margins below are rational inequalities; the
supplementary script \path{check_two_core_upper.py} checks them using
exact fractions.

\begin{lemma}[The zero threshold]\label{t:zero}
If at least three coefficients are nonzero, then
$\Pp(X=0)\le3/8$. Equality forces four equal nonzero coefficients.
\end{lemma}
\begin{proof}
Absorb signs and remove zeros. Condition on all but three coefficients.
The signs of three positive coefficients giving any fixed value form
an antichain, so have probability at most $3/8$. For completeness,
a random maximal chain of the three-dimensional Boolean lattice meets
an antichain at most once, and each antichain member is met with
probability at least $1/3$, proving the count bound three.

For three positive weights $u\ge v\ge w$, an atom has multiplicity
three only if $u=v=w$, and its value is then $\pm u$.
Indeed the extreme values $\pm(u+v+w)$ are singletons; the remaining
values are $\pm(u+v-w),\pm(u-v+w),\pm(-u+v+w)$.
Three coincide only when the three positive expressions coincide,
which gives $u=v=w$; if $-u+v+w\le0$, a triple equality would
force $v=0$ or $w=0$.
Thus equality in the conditional bound forces the three selected
weights to equal some $b>0$ and the residual $R$ to be supported on
$\{-b,b\}$. Its largest outcome is its coefficient sum $\ell=b$,
while $\E R^2=b^2$. Hence
$\ell^2=\sum_{i\in R}b_i^2$, forcing exactly one nonzero residual
coefficient. Normalization gives four coefficients $b=1/2$.
\end{proof}

\begin{lemma}[Translated residual intervals]\label{t:translated}
If a finite Rademacher sum $R$ has variance $V>0$ and every coefficient
has magnitude at most $x>0$, then every interval of radius $x$ has
probability less than $(8/5)x/\sqrt V$.
\end{lemma}
\begin{proof}
After normalization, Keller--Klein bounds the difference of each CDF
(and its left limits) from the Gaussian by $c(x/\sqrt V)$. The Gaussian
mass of an interval of radius $x/\sqrt V$ is at most $2c(x/\sqrt V)$.
Subtracting the two CDF values therefore gives at most $4c(x/\sqrt V)$,
which is strictly less than $(8/5)x/\sqrt V$.
\end{proof}

\begin{lemma}[Exact two-core radius]\label{t:radius}
Suppose $a\le1/\sqrt2$, $x>0$, and
\[
 U_X(x)\ge T:=\max\{3/16,1/4-c(r)\},\qquad r=\sqrt{1-2a^2}.
\]
Then $x\ge r$ and $\Pp(|X|\le x)\le1/2$.
\end{lemma}
\begin{proof}
For $1/\sqrt7\le x<1/\sqrt3$, the strict tail is at least $3/8$,
so $U_X(x)\le5/16-c(1/\sqrt7)<11/64<3/16$.
For $1/\sqrt3\le x<1$, it is at least $7/32$, giving
$U_X(x)\le25/64-c(1/\sqrt3)<3/16$.
For $x\ge1$, $U_X(x)\le1-\Phi(1)<3/16$.
Consequently $x<1/\sqrt7$, where the first tail phase gives
$p:=\Pp(|X|\le x)\le1/2$. Thus $U_X(x)\ge T\ge3/16$
implies $x\le q<4/25$.

Also $p=2U_X(x)+2c(x)>3/8$. If three coefficients were greater
than $x$, conditioning on the others and applying the same antichain
argument would give $p\le3/8$. Hence, after sorting, $b_3\le x$.
If there are at most two nonzero coefficients, normalization and
$a\le1/\sqrt2$ force the equal pair and $r=0$, proving the claim.
Otherwise put $V=\sum_{i\ge3}b_i^2>0$, $s=b_1+b_2$, $h=b_1-b_2$.
Lemma~\ref{t:translated}, applied after conditioning the first two
signs, gives
\[
 V<\left(\frac{(8/5)x}{p}\right)^2.
\]
The law of $X$ is the equal mixture of the Rademacher laws
$h\varepsilon+R$ and $s\varepsilon+R$, with the displayed sign
independent of $R$. The former has variance
\[
 h^2+V=1-2b_1b_2\ge r^2.
\]

Assume for contradiction that $x<r$, and put
$u=\min\{r,q\}$ and $t=x/u$. Then $0<t\le1$; if $t=1$,
we necessarily have $u=q<r$. Since $T=1/4-c(u)$,
\[
 p\ge\frac12-2(c(u)-c(x))
 \ge\frac{93+32t}{250}=:P(t).
\]
Here $u\le4/25$ and $c'\le2/5$ were used. In particular
\[
 V<v(t):=\left(\frac{64t}{93+32t}\right)^2,
 \qquad x\le\frac{4t}{25}.
\]
The core sum satisfies $s\ge(1-V)/a\ge\sqrt2(1-V)$, so
\[
 s-x>\frac{141}{100}(1-v(t))-\frac{4t}{25}=:d(t)>0.
\]
The positivity follows from $v(t)\le v(1)=4096/15625$ and
$d(1)=1375589/1562500>0$.

The balanced component has central probability at most $1-f(t)$,
using the appropriate phase above. In the last phase the same bound
$1-7/32$ applies also when $t=1$, because its actual normalized
threshold is strictly less than one by $x<r$.
For the other component, symmetry and one-sided Chebyshev give
\[
 \Pp(|s\varepsilon+R|\le x)
 =\Pp(s-x\le R\le s+x)
 \le\frac{V}{2(s-x)^2}.
\]
Therefore
\[
 p<\frac{1-f(t)}2+F(t),\qquad
 F(t):=\frac{v(t)}{4d(t)^2}.
\]
The function $v$ increases and $d$ decreases on $[0,1]$, so $F$
increases. In each phase, replace $t$ in $F$ by the rational upper
endpoint in the next table, and replace $t$ in $P$ by the rational
lower endpoint. Each displayed positive margin is an exact rational
lower bound for $P(t)-(1-f(t))/2-F(t)$.
\[
\begin{array}{c|c|c|c}
(1-f(t))/2&t\text{ lower}&t\text{ upper}&\text{margin greater than}\\\hline
1/4&0&38/100&11/100\\
35/128&37/100&45/100&13/100\\
5/16&44/100&58/100&9/100\\
3/8&57/100&82/100&2/100\\
25/64&81/100&1&49/100000
\end{array}
\]
The square-root phase endpoints lie inside the listed rational bounds,
as is verified by squaring. Thus the upper bound for $p$ is strictly
below $P(t)$ in every phase, a contradiction. Hence $x\ge r$.
\end{proof}

\begin{proof}[Proof of Theorem~\ref{t:main}]
For $a<1/\sqrt2$, there must be at least three nonzero coefficients,
so Lemma~\ref{t:zero} gives $U_X(0)\le3/16$, with equality only
for four-flat. For positive thresholds competitive with $T$,
Lemma~\ref{t:radius} gives
\[
 U_X(x)\le1/4-c(x)\le1/4-c(r)\le T.
\]
The endpoint $a=1/\sqrt2$ also permits the equal pair, whose value
is $1/4$ at zero. For any larger cap, the same value remains maximal:
the localization estimates above give $U_X(x)<3/16$ for
$x\ge1/\sqrt7$, while the first strict-tail phase gives
$U_X(x)\le1/4-c(x)<1/4$ for $0<x<1/\sqrt7$. At zero,
conditioning on one nonzero sign gives an atom of mass at most
$1/2$. Equality requires exactly two equal nonzero coefficients,
because three or more are covered by Lemma~\ref{t:zero}.

Four-flat is admissible for every $a\ge1/2$ and gives $3/16$.
Whenever $1/4-c(r)\ge3/16$, one has $r\le q<4/25<a$.
Thus $(a,a,r)$ is admissible and its central mass at $x=r$ is
$1/2$; it gives the other branch. The branch is strictly increasing
in $a$, yielding the unique transition $a_0$.

It remains to classify positive-threshold equality. It requires
$x=r$ and $p=1/2$. The two-coefficient endpoint is settled, so assume
$V>0$. By the preceding conditioning argument, $b_3\le x$ and
\[
 V<\bigl((8/5)x/(1/2)\bigr)^2
 <4096/15625.
\]
Consequently
$b_2^2=1-V-b_1^2>1/2-4096/15625>(4/25)^2>x^2$.
For each residual outcome $z$, an interval of radius $x$ can contain
at most two of the four values $-s,-h,h,s$; if it contains two,
they must be the central pair $-h,h$, because $s-h=2b_2>2x$.
The average conditional central probability equals $1/2$, so each
conditional probability equals $1/2$. Therefore $|z|+h\le r$ for
every residual outcome. Taking $\ell=\sum_{i\ge3}b_i$ gives
\[
 r^2\le V\le\ell^2\le(r-h)^2\le r^2.
\]
Equality throughout forces $h=0$, $b_1=b_2=a$, and one nonzero
residual coefficient $r$. Together with Lemma~\ref{t:zero}, this
proves all equality and tie cases.
\end{proof}

\paragraph{A general balanced-core observation.}
The decisive variance inequality extends to every even core size $M$.
Generate the uniform balanced sign slice by a uniform perfect matching,
then independently orient each matched pair with opposite signs.
Conditional on a matching, the balanced core plus residual is a
Rademacher sum of variance
\[
 V+\sum_{\{i,j\}\text{ matched}}(b_i-b_j)^2
 =1-2\sum_{\{i,j\}\text{ matched}}b_i b_j\ge1-Ma^2.
\]
Thus the same Hollom--Portier bound applies to every matching.
For $M\ge4$, the missing global step is threshold localization and
control of the other core layers. No all-cap upper formula is inferred
from this observation alone.

\section{Lower discrepancy: two finite certificates}\label{sec:lower-certificates}
Write $Q(t)=1-\Phi(t)$ and
\[
 L_X(x)=c(x)-\tfrac12\Pp(|X|<x)=\Pp(X\ge x)-Q(x),\qquad x>0.
\]
The strict central interval is essential. Besides the Keller--Klein
comparison and the Hollom--Portier profile $f$ in
\eqref{eq:HP-input}, this section uses Tomaszewski's theorem
$\Pp(|Z|\le1)\ge1/2$ for a normalized Rademacher sum~\cite{KK}.
All finite certificates below concern explicit scalar inequalities;
their use of validated interval arithmetic is separate from Lean.

\begin{theorem}[The five-flat lower plateau]\label{thm:five-flat-lower}
If $\max_i|b_i|\le1/2$, then
\[
 \sup_{x>0}L_{X_b}(x)\le c(1/\sqrt5).
\]
Equality holds only for the five-flat vector, at $x=1/\sqrt5$.
Consequently the lower-discrepancy cap optimum equals $c(1/\sqrt5)$
for every $a\in[1/\sqrt5,1/2]$.
\end{theorem}

\begin{theorem}[Exact lower discrepancy in the first missing band]\label{thm:three-core-lower}
For $1/\sqrt5\le a\le1/\sqrt3$, put
\[
 G_3(a)=\max\{1/\sqrt5,a-\sqrt{1-3a^2}\}=\mathfrak G(a).
\]
Then
\begin{equation}\label{eq:lower-exact-first-band}
 \sup_{\substack{N<\infty,\ \sum_i b_i^2=1\\ |b_i|\le a}}\ \sup_{x>0}L_{X_b}(x)=c(G_3(a)).
\end{equation}
Equality occurs precisely at $x=G_3(a)$ for a maximizing support-gap
vector: five-flat on the flat branch, $(a,a,a,\sqrt{1-3a^2})$ on
the saturated branch, and both at the unique transition.
Together with the earlier high-cap result~\cite{Prior}, the identity
$\mathfrak L(a)=c(\mathfrak G(a))$ and this equivalence of lower-
discrepancy and support-gap optimizers now hold for every
$a\ge1/\sqrt5$.

The positive-discrepancy theorem consequently gives
\begin{equation}\label{eq:full-half-third}
 \mathfrak K(a)=\max\{3/16,c(a-\sqrt{1-3a^2})\},
 \qquad 1/2\le a\le1/\sqrt3.
\end{equation}
Put $h=c^{-1}(3/16)$ and
$a_3=(h+\sqrt{4-3h^2})/4$. Below $a_3$ the unique maximizing type
is four-flat, at positive-discrepancy threshold zero; above $a_3$
it is $(a,a,a,\sqrt{1-3a^2})$, at lower-discrepancy threshold
$a-\sqrt{1-3a^2}$. At $a_3$ precisely both types maximize.
Permutations, coefficient sign changes and zero padding are immaterial.
\end{theorem}

\subsection{The four-shift envelope}
Absorb coefficient signs and order $m=b_1\ge b_2\ge b_3\ge\cdots$.
Write $R=\sum_{i\ge4}b_i\varepsilon_i$ and
$\sigma^2=1-m^2-b_2^2-b_3^2>0$. Set
\[
 \rho=\min\{1,b_3/\sigma\},\quad C=c(\rho),\qquad
 (q_0,q_1,q_2,q_3)
 =(m+b_2+b_3,\,|b_2+b_3-m|,\,m-b_2+b_3,\,m+b_2-b_3).
\]
The distribution of $X$ is the uniform mixture of the four symmetric
laws $q_j\varepsilon+R$. For $t>0$ define
\begin{equation}\label{eq:closed-tail-envelope}
 T_C(t)=\min\left\{\frac12,\frac1{2t^2},e^{-t^2/2},Q(t)+C,
 \quad\frac14\ \text{if }t>1\right\}.
\end{equation}
An inapplicable last entry is omitted. For $u\in\mathbb R$ set
\begin{equation}\label{eq:open-tail-envelope}
 A_C(u)=
 \begin{cases}
 \max\{1-T_C(-u),Q(u)-C,0\},&u<0,\\
 \max\{f(u)/2,Q(u)-C,0\},&u\ge0,
 \end{cases}
\end{equation}
where $f(u)=0$ for $u\ge1$ and otherwise has the stated strict-tail
values, with the smaller value at every jump. For each $j$, put
\[
 u_j=(q_j-x)/\sigma,\quad v_j=(q_j+x)/\sigma,
 \qquad
 H_j=\max\{0,A_C(u_j)-T_C(v_j),Q(u_j)-Q(v_j)-2C\}.
\]
Additional valid lower bounds on $\Pp(R/\sigma>u_j)$ may be included
in the maximum defining $A_C(u_j)$. We use only the explicit
additional bounds described below. The resulting scalar envelope is
\begin{equation}\label{eq:four-shift-envelope}
 E(m,b_2,b_3,x)=\min\left\{c(m),\,
 c(x)-\frac18\sum_{j=0}^3H_j,\,
 \frac1{2x^2}-Q(x)\right\}.
\end{equation}

\begin{lemma}\label{lem:four-shift-valid}
With these definitions, $L_X(x)\le E(m,b_2,b_3,x)$ for every $x>0$.
\end{lemma}
\begin{proof}
For the normalized residual $Z=R/\sigma$, symmetry and Chebyshev give
\[
 \Pp(Z\ge t)\le\min\{1/2,1/(2t^2)\}.
\]
The exponential bound follows
from $\E e^{sZ}=\prod_i\cosh(sb_i/\sigma)\le e^{s^2/2}$ and Markov,
optimized at $s=t$. Keller--Klein, including CDF left limits, gives
$\Pp(Z\ge t)\le Q(t)+C$ and $\Pp(Z>u)\ge Q(u)-C$.
For $t>1$, Tomaszewski gives
$\Pp(Z\ge t)\le\Pp(Z>1)\le1/4$. These prove
\eqref{eq:closed-tail-envelope}. If $u<0$, symmetry gives
$\Pp(Z>u)=1-\Pp(Z\ge-u)$; if $u\ge0$, Hollom--Portier gives
$\Pp(Z>u)=\frac12\Pp(|Z|>u)\ge f(u)/2$.
Thus \eqref{eq:open-tail-envelope} is valid, with the stated endpoints.
Finally,
\[
 \Pp(|q_j\varepsilon+R|<x)
 =\Pp(q_j-x<R<q_j+x)
 =\Pp(Z>u_j)-\Pp(Z\ge v_j)\ge H_j.
\]
Averaging the four laws proves the middle bound in
\eqref{eq:four-shift-envelope}; the other two are the global
Keller--Klein bound and one-sided Chebyshev for $X$.
\end{proof}

\subsection{Certified scalar inequalities}\label{subsec:scalar-interface}
The following is the exact interface between the analysis and the
computation. Every box coordinate and every bisection is rational.
Only the evaluation of square roots, exponentials and Gaussian integrals
uses validated real intervals.

\begin{lemma}[Finite scalar certificates]\label{lem:three-certificates}
The envelope \eqref{eq:four-shift-envelope} has the following strict
bounds, on the full indicated domains.
\begin{enumerate}
\item If $1/2\le m\le1/\sqrt3$, $0\le b_3\le b_2\le m$,
$447/1000\le x\le9/5$, $\sigma>0$, and $x>G_3(m)$, then
$E<c(G_3(m))$,
where the extra bounds $A_C(u_1),A_C(u_2)\ge7/64$ are included,
and $A_C(u_3)\ge7/64$ is included whenever $b_3\ge\sqrt{1-3m^2}$.
\item If $1/\sqrt5\le m\le1/2$, $0\le b_3\le b_2\le m$,
$447/1000\le x\le9/5$, and $q_1<x$, then $E<c(1/\sqrt5)$,
where the extra bound $A_C(u_1)\ge1/2$ is included.
\item If $1/\sqrt5\le m\le1/2$, $0\le b_3\le b_2\le m/8$,
and $447/1000\le x\le9/5$, then $E<c(1/\sqrt5)$ without
extra bounds.
\end{enumerate}
\end{lemma}
\begin{proof}[Proof and reproducibility of the scalar certificates]
The respective partitions have $13666$, $14303$ and $3$ leaves.
They are supplied, compressed with \texttt{xz}, as
\path{lower_exact_gap_pilot.json.xz},
\path{lower_five_flat_pilot_B.json.xz}, and
\path{lower_five_flat_pilot_C.json.xz}.
Their names record their origin in a nonrigorous partition search;
none of the floating point values in those files is trusted by the
verifiers. The verifiers read only the finite bisection paths, reconstruct
every endpoint with rational arithmetic, verify complete coverage, and
then prove a strictly positive interval margin on every leaf.

Here are complete details of the enclosure algorithm.
Parameterize $w=m^2$, $b_2=m\alpha$, $b_3=m\alpha\beta$.
In case (i),
\[
 (w,\alpha,\beta,x)\in[1/4,1/3]\times[0,1]^2
 \times[447/1000,9/5];
\]
in cases (ii),(iii), the $w$ interval is $[1/5,1/4]$, the $x$
interval is $[447/1000,9/5]$, and the $\alpha$ interval is respectively
$[0,1]$ or $[0,1/8]$. A digit $2k+e$ in a leaf path means bisect
coordinate $k\in\{0,1,2,3\}$ and choose half $e\in\{0,1\}$.
The exact prefix-tree check requires that every nonleaf has precisely
the two children $2k,2k+1$ for one common $k$; it also checks absence
of duplicate or prefix leaves and exact total normalized volume one.
Thus no region can be omitted.

For a box write subscripts $-$ and $+$ for its coordinate endpoints,
and $m_\pm=\sqrt{w_\pm}$. Valid residual bounds are
\[
 V_-=1-w_+(1+\alpha_+^2(1+\beta_+^2)),\qquad
 V_+=1-w_-(1+\alpha_-^2(1+\beta_-^2)).
\]
Both are nonnegative on the specified domains, and $V_+>0$.
Set $s_\pm=\sqrt{V_\pm}$ and
$C_+=c(\min\{1,m_+\alpha_+\beta_+/s_-\})$, taking $C_+=c(1)$
if $s_-=0$. The elementary interval bounds for the four shifts are
\[
\begin{array}{c|cc}
j&q_{j,-}&q_{j,+}\\\hline
0&m_-(1+\alpha_-(1+\beta_-))&m_+(1+\alpha_+(1+\beta_+))\\
2&m_-(1-\alpha_+(1-\beta_-))&m_+(1-\alpha_-(1-\beta_+))\\
3&m_-(1+\alpha_-(1-\beta_+))&m_+(1+\alpha_+(1-\beta_-)).
\end{array}
\]
For $j=1$, take the minimum and maximum of the absolute value on
$[\alpha_-(1+\beta_-)-1,\alpha_+(1+\beta_+)-1]$, then multiply
by $m_-$ and $m_+$ respectively. In case (i) replace $x_-$ by
$\max\{x_-,1/\sqrt5,m_- -\sqrt{1-3w_-}\}$; this only restricts
to a larger set than the required condition $x>G_3(m)$.
The target is bounded below by $c(G_3(m_-))$ using monotonicity.

An upper bound for $u_j$ is
$(q_{j,+}-x_-)/s_-$ when its numerator is nonnegative, and
$(q_{j,+}-x_-)/s_+$ when negative; use $+\infty$ if the first
case has $s_-=0$. A lower bound for $v_j$ is
$(q_{j,-}+x_-)/s_+>0$. Apply the formulas
\eqref{eq:closed-tail-envelope}--\eqref{eq:four-shift-envelope}
to these endpoints, using $C_+$, $c(m_+)$, $c(x_+)$ and
$1/(2x_-^2)-Q(x_+)$. Monotonicity gives an upper bound for $E$:
$T_C$ decreases with its argument and increases with $C$, while
$A_C$ decreases with its argument and with $C$.
For extra bounds, cases (i),(ii) use exactly their stipulated entries;
the optional $j=3$ entry in (i) is used only when the exact rational
comparison $w_-\alpha_-^2\beta_-^2\ge1-3w_-$ holds.
If an interval comparison cannot determine a sign, the verifier uses
a weaker bound, never the unproved branch. At a phase boundary it
uses only a phase whose endpoint is proved strictly larger. The
exceptional forced $7/64$ entries are justified by a strict analytic
inequality below, not by rounding.

The programs \path{verify_lower_exact_gap_arb.py} and
\path{verify_lower_five_flat_arb.py} implement these formulas with
\texttt{Fraction} for exact rational data and \texttt{python-flint}
0.9.0 (Arb) for real enclosures~\cite{Arb,PythonFlint}. A floating midpoint is used only to
choose one among already valid upper or lower bounds; correctness
never relies on that choice selecting the best bound. The acceptance
test is an interval comparison proving the entire margin strictly
positive. All $27972$ leaves pass at both $128$ and $192$ bits.
The supplied JSON reports contain each leaf margin, partition and
verifier SHA-256 hashes, coverage checks, precision and execution
metadata. The commands are
\begin{verbatim}
python verify_lower_exact_gap_arb.py 192
python verify_lower_five_flat_arb.py B 192
python verify_lower_five_flat_arb.py C 192
\end{verbatim}
This proves the stated scalar inequalities with validated arithmetic.
It is not a Lean kernel check or a replay of the external Hollom--Portier
computation.
\end{proof}

\subsection{Proofs of the probability conclusions}
\begin{proof}[Proof of Theorem~\ref{thm:three-core-lower}]
The five-flat plateau covers actual maxima $m\le1/2$; its only equality
is five-flat at $x=1/\sqrt5$. Assume $1/2<m\le a\le1/\sqrt3$.
For $x\le G_3(m)$, trivially $L_X(x)\le c(x)\le c(G_3(m))$.
We show that $x>G_3(m)$ gives strict inequality.
If $x\ge9/5$, one-sided Chebyshev gives
$L_X(x)\le25/162<c(1/\sqrt5)\le c(G_3(m))$.
Otherwise $x>G_3(m)\ge1/\sqrt5>447/1000$, so the variables lie
in the domain of Lemma~\ref{lem:three-certificates}(i).
If $\sigma=0$, normalization and $m\le1/\sqrt3$ force the
three-flat vector. For $x>G_3(m)=1/\sqrt3$, its strict central
probability is at least $3/4$, hence $L_X(x)\le1/8<c(1/\sqrt5)$.

To justify the additional tail entries in (i), set
$r_3=\sqrt{1-3m^2}$. Since $q_1,q_2\le m$ and $\sigma\ge r_3$,
the strict inequality $x>G_3(m)\ge m-r_3$ gives
$q_1-x<\sigma$ and $q_2-x<\sigma$. Thus $u_1,u_2<1$;
Hollom--Portier gives $A_C(u_1),A_C(u_2)\ge7/64$ if these
arguments are nonnegative, and symmetry gives at least $1/2$ otherwise.
For $j=3$, let $d=b_2-b_3\ge0$. If $b_3\ge r_3$, then
\[
 \sigma^2-r_3^2=2m^2-b_2^2-b_3^2
 \ge d^2+2b_3d\ge d^2+2r_3d.
\]
The first inequality follows by subtraction, leaving
$2m^2-2b_2^2\ge0$. Thus $\sigma\ge r_3+d$, and
$q_3-x=m+d-x<r_3+d\le\sigma$. This proves the optional last
entry. The scalar certificate and Lemma~\ref{lem:four-shift-valid}
now give the required strict bound.

The function $G_3$ is nondecreasing. Therefore every admissible law
has lower discrepancy at most $c(G_3(a))$, and equality must have
$x=G_3(a)$ and an empty strict central interval. Equivalently,
$g(b)=\mathfrak G(a)$, so the support-gap theorem proves both
attainment and the complete optimizer classification. This establishes
\eqref{eq:lower-exact-first-band}. The earlier theorem for
$a\ge1/\sqrt3$ gives the claimed combined range.

On $[1/2,1/\sqrt3]$, the positive optimum is $3/16$ and
$c(1/\sqrt5)<3/16$. Taking the maximum of the positive and lower
optima proves \eqref{eq:full-half-third}, with the equality cases
inherited from their separate classifications. Solving
$a-\sqrt{1-3a^2}=h$ gives the displayed transition $a_3$.
\end{proof}

\begin{proof}[Proof of Theorem~\ref{thm:five-flat-lower}]
Write $\lambda=1/\sqrt5$ and $H=c(\lambda)$.
If $m<\lambda$, Keller--Klein gives $L_X\le c(m)<H$.
If $x<\lambda$, the trivial bound $L_X(x)\le c(x)$ is strict.
At $x=\lambda$, equality means $\Pp(|X|<\lambda)=0$, hence
$g(b)\ge\lambda$. The exact support-gap theorem gives
$\mathfrak G(1/2)=\lambda$ with only five-flat attaining it,
so this case has precisely the asserted equality type. If $x\ge9/5$, then
$L_X(x)\le25/162<H$, a scalar comparison in the certificates.
We may therefore assume $\lambda\le m\le1/2$ and
$\lambda<x<9/5$. Here $\sigma^2\ge1-3m^2\ge1/4$,
so the four-shift envelope is always defined.

If $q_1<x$, symmetry gives
$\Pp(Z>u_1)\ge\Pp(Z\ge0)\ge1/2$.
Lemma~\ref{lem:three-certificates}(ii) gives $L_X(x)<H$.
If $q_1\ge x$ and $b_2+b_3<m$, then
$b_2+b_3\le m-x\le m-\lambda$ and hence
\[
 b_2/m\le1-\lambda/m\le1-2\lambda<1/8.
\]
Lemma~\ref{lem:three-certificates}(iii) again gives $L_X(x)<H$.

It remains to treat the triangle case $b_2+b_3\ge m$ and $q_1\ge x$.
Then $b_2+b_3\ge m+x$, so
\[
 x\le q_1\le q_2\le q_3,\qquad x\le m\le1/2,\qquad
 q_3-x\le2(m-x)\le1-2\lambda<\sigma/\sqrt7.
\]
The last comparison follows from $\sigma\ge1/2$ and
$2-4\lambda<1/\sqrt7$, also verified in the certificates.
For $j=1,2,3$, the first Hollom--Portier phase and one-sided Chebyshev
therefore give
\[
 \Pp(q_j-x<R<q_j+x)
 \ge\frac14-\frac{\sigma^2}{2(q_j+x)^2}
 \ge\frac14-\frac{\sigma^2}{8x^2}.
\]
Also $b_2^2+b_3^2\ge(b_2+b_3)^2/2\ge(m+x)^2/2$, whence
$\sigma^2\le1-m^2-(m+x)^2/2\le1-3x^2$, with equality
in the last inequality only if $m=x$. Therefore
\[
 L_X(x)\le c(x)-\frac3{32}+\frac{3\sigma^2}{64x^2}
 \le c(x)-\frac{15}{64}+\frac3{64x^2}=:J(x).
\]
On $[\lambda,1/2]$,
$J'(x)=\varphi(x)-3/(32x^3)<2/5-3/4<0$, while
$J(\lambda)=c(\lambda)$ exactly. Thus $L_X(x)<H$ when
$x>\lambda$. This exhausts every remaining case and proves both the bound and
its equality assertion.
\end{proof}

\section{The exact full discrepancy down to the five-flat cap}\label{sec:four-core-full}
The preceding lower certificate combines with a four-coordinate
localization argument to determine the full discrepancy throughout the
next cap interval. No additional numerical optimization is used here.

\begin{theorem}\label{thm:full-five-half}
Let $\lambda=1/\sqrt5$, $H=c(\lambda)$, and
$r_4(a)=\sqrt{1-4a^2}$. Then
\begin{equation}\label{eq:full-five-half}
 \mathfrak K(a)=\max\{H,\,3/16-c(r_4(a))\},
 \qquad \lambda\le a\le1/2.
\end{equation}
Put $q_4=c^{-1}(3/16-H)$ and $a_4=\frac12\sqrt{1-q_4^2}$.
For $\lambda\le a<a_4$, the unique maximizing type is five-flat,
at lower-discrepancy threshold $\lambda$. For $a_4<a\le1/2$,
it is $(a,a,a,a,r_4(a))$, at positive-discrepancy threshold $r_4(a)$.
At $a_4$, precisely both types maximize. At $a=1/2$ the latter
vector is four-flat and its threshold is zero.
\end{theorem}

\begin{lemma}[Four-coordinate localization]\label{lem:four-core-localization}
Suppose $\lambda\le a\le1/2$, $x>0$, and $U_X(x)\ge H$.
Then exactly four coefficients are greater than $x$,
$\Pp(|X|\le x)\le3/8$, and $x\ge r_4(a)$.
\end{lemma}
\begin{proof}
The scalar estimates needed below are
\[
 H>11/64>17/100,\quad c(1/\sqrt7)>9/64,\quad
 c(1/\sqrt3)>43/200,\quad c(4/5)>7/25,\quad Q(1)<17/100,
\]
\[
 c(1/5)>5/64,\qquad c(1/25)>1/64,\qquad
 c(x)>39x/100\quad(0<x\le1/5).
\]
They follow from $199/500<\varphi(0)<2/5$,
$e^{-t^2/2}\ge1-t^2/2$, the elementary square-root bounds
$1/\sqrt5>447/1000$, $1/\sqrt7>377/1000$,
$1/\sqrt3>577/1000$, and the exact rational comparisons in
\path{check_four_core_upper.py}. For example,
$H>(199/500)(29/30)(447/1000)>11/64$, and on $[0,1/5]$
we have $\varphi(x)>(199/500)(49/50)>39/100$.

The strict-tail profile first forces $x<1/\sqrt7$. Indeed, in
$[1/\sqrt7,1/\sqrt3)$ it gives
$U_X(x)\le5/16-c(1/\sqrt7)<11/64<H$; in
$[1/\sqrt3,\sqrt{2/3})$ it gives
$U_X(x)\le3/8-c(1/\sqrt3)<4/25<H$; and in
$[\sqrt{2/3},1)$ it gives
$U_X(x)\le25/64-c(4/5)<17/100<H$. Finally $x\ge1$
gives $U_X(x)\le Q(1)<H$.
For $x<1/\sqrt7$ the first tail phase gives
$p:=\Pp(|X|\le x)\le1/2$. Consequently
$c(x)\le1/4-H<5/64$, so $x<1/5$, and
\begin{equation}\label{eq:four-core-p-lower}
 p=2U_X(x)+2c(x)>17/50+(39/50)x.
\end{equation}

If at most two coefficients exceed $x$, condition on the first two
signs. The remaining variance is at least $1-2a^2\ge1/2$, and
all remaining coefficients are at most $x$. The translated interval
bound gives
$p<(8\sqrt2/5)x<(12/5)x$. This contradicts
\eqref{eq:four-core-p-lower}, since
$(12/5-39/50)(1/5)=81/250<17/50$.
Thus at least three coefficients exceed $x$. The sign patterns of
any three such coefficients falling in an interval of length $2x$
form an antichain, so $p\le3/8$ by Sperner's bound.
It follows that $c(x)\le3/16-H<1/64$, hence $x<1/25$.

Exactly three coefficients cannot exceed $x$: conditioning on them
leaves variance at least $1-3a^2\ge1/4$, and the translated interval
bound gives $p<(16/5)x<16/125<17/50$.
Nor can five coefficients exceed $x$: conditioning on any five gives
the antichain bound $p\le\binom52/2^5=5/16<17/50$.
Thus exactly four exceed $x$.

Sort them as $b_1\ge b_2\ge b_3\ge b_4>x$, and let
$R=\sum_{i\ge5}b_i\varepsilon_i$ have variance $V$.
If $V=0$, then normalization and $a\le1/2$ force four-flat,
and $r_4(a)=0$, proving the radius assertion.
Otherwise, the translated interval bound and
$p>17/50$, $x<1/25$ give
\begin{equation}\label{eq:four-core-small-V}
 V<\left(\frac{(8/5)(1/25)}{17/50}\right)^2
 =\left(\frac{16}{85}\right)^2<9/250.
\end{equation}
Suppose $x<r_4(a)$. The six balanced sign patterns have total weight
$6/16$. Their uniform law can be generated by choosing one of the
three perfect matchings of four indices uniformly, and independently
orienting each matched pair with opposite signs. Every balanced sign
pattern is then counted twice among the twelve equally likely
outcomes. Conditional on a matching, the balanced core plus $R$ is
a Rademacher sum of variance
\[
 V+\sum_{\{i,j\}\text{ matched}}(b_i-b_j)^2
 =1-2\sum_{\{i,j\}\text{ matched}}b_i b_j
 \ge1-4a^2=r_4(a)^2.
\]
Its normalized threshold is strictly below one. Thus the last
Hollom--Portier phase, or any stronger preceding phase, bounds its
central probability by $25/32$.

For every nonbalanced sign pattern the core magnitude is at least
$S-2a$, where $S=\sum_{i=1}^4b_i$. Since
$S\ge(1-V)/a$, we have
\[
 S-2a\ge2a-V/a>
 \frac{894}{1000}-\frac{9/250}{447/1000}>
 \frac{813}{1000}.
\]
Here $a\le1/2$ gives $1/a-2a\ge2a$, and $a\ge\lambda>447/1000$.
For a nonbalanced magnitude $s$, symmetry and one-sided Chebyshev
therefore give
\[
 \Pp(|s\varepsilon+R|\le x)
 =\Pp(s-x\le R\le s+x)
 \le\frac{V}{2(s-x)^2}
 <\frac{9/250}{2(773/1000)^2}<\frac{31}{1000}.
\]
Combining the six balanced and ten nonbalanced patterns yields
\[
 p<\frac6{16}\frac{25}{32}+\frac{10}{16}\frac{31}{1000}
 <\frac{17}{50},
\]
contrary to \eqref{eq:four-core-p-lower}. Hence $x\ge r_4(a)$.
\end{proof}

\begin{proof}[Proof of Theorem~\ref{thm:full-five-half}]
The lower discrepancy is at most $H$ by
Theorem~\ref{thm:five-flat-lower}, with only five-flat attaining it.
For positive $x$ competitive with $H$, the localization lemma gives
\[
 U_X(x)\le3/16-c(x)\le3/16-c(r_4(a)).
\]
At zero, if $a<1/2$ there are at least five nonzero coefficients;
the five-sign atom bound gives $U_X(0)\le5/32<H$. If $a=1/2$,
then $U_X(0)\le3/16$, with equality only for four-flat.
This proves the upper bound in \eqref{eq:full-five-half}.

Five-flat is admissible throughout the interval and attains $H$.
The other branch can compete only when
$c(r_4(a))\le3/16-H<1/64$, which implies
$r_4(a)<1/25<a$. The vector $(a,a,a,a,r_4(a))$ is therefore
admissible; precisely its six balanced core patterns fall in the
closed central interval of radius $r_4(a)$, giving
$U=3/16-c(r_4(a))$. This branch increases strictly with $a$,
so the displayed transition is unique.

For positive-threshold equality on that branch, the localization
lemma forces $x=r_4(a)$, exactly four coefficients greater than $x$,
and $p=6/16$. For each residual outcome, the qualifying four-sign
patterns form an antichain of size at most six. An antichain of size
six in the four-dimensional Boolean lattice is exactly the middle
layer: the maximal-chain identity
$\sum_{A\in\mathcal A}\binom4{|A|}^{-1}\le1$
forces every member to have size two. Since the average conditional
probability is $6/16$, every residual outcome must include all six
balanced patterns. Their extreme values are $\pm d$, where
$d=b_1+b_2-b_3-b_4\ge0$. Writing
$\ell=\sum_{i\ge5}b_i$ gives $\ell+d\le r_4(a)$, and hence
\[
 r_4(a)^2\le V\le\ell^2\le(r_4(a)-d)^2\le r_4(a)^2.
\]
All inequalities are equalities. The core energy is $4a^2$, so all
four core coefficients equal $a$; the tail satisfies $\ell^2=V$,
so it has one nonzero coefficient, of magnitude $r_4(a)$.
The zero-threshold endpoint and the lower equality classification
already proved give exactly the stated optimizer and tie cases.
\end{proof}

\section{Verification scope}\label{sec:verification}
The written mathematical proofs, external theorem inputs, exact-arithmetic
checks, and Lean kernel checks are distinct parts of the verification
record. Successful compilation of a selected module does not certify
unformalized analytic statements elsewhere in the paper.

\paragraph{Kernel-checked scope.}
The inherited support-gap development proves the exact all-dimensional
value, actual attainment, and optimizer classification from finite
signed sums. The new module \code{GlobalGapRigidity.lean} proves the
sharp global inverse and its equality classification on the original
coefficient indices. It derives the actual core from the given vector;
the majority structure, normalized variance identities, and finite-slice
inequality are proved rather than assumed. The wrappers
\code{signed_supportGap_inverse} and
\code{signed_supportGap_inverse_equality_iff} accept arbitrary signed
coefficients and their actual minimum absolute signed sum.

The wrapper in \code{ActualGlobalGapStability.lean} derives the odd cap band and actual
core, constructs a normalized comparison vector obeying the cap, proves
the sharp $\sqrt{2+\sqrt2}$ distance bound, and proves that the comparison
is an actual optimizer under the stated phase separation or at a tie.
The comparison is given on the original indices with one zero coordinate
added when necessary. This is a concrete finite-dimensional distance
statement; sorting magnitudes gives the distance convention used here.

\code{FixedFlatTradeoff.lean} proves both displayed fixed-flat
inequalities, including their exact radical form, from the actual input
vector. The equality-shape analysis of
Proposition~\ref{prop:fixed-flat-gap} is supplied by its written proof.
\code{ExactSaturatedModulus.lean} proves the exact nonlinear saturated
modulus and its complete equality equivalence: every core coordinate
equals $a$, the largest tail is $t_+$, and the remaining tail has at most
one nonzero coefficient with sum $t_-$. The module
\code{ActualSaturatedModulus.lean} joins the bound to arbitrary signed
input, the actual support gap, zero padding, and global phase selection.
The separate \code{ExactModulusAttainment.lean} realizes the entire
two-residual equality family for every odd core size and every feasible
local deficit, including both endpoints. It proves the actual gap, cap,
normalization, core size, and distance from explicit coefficients.

The examples in \code{GapRigidityExamples.lean} prove the actual
saturated equality family and the normalized endpoint counterexample
$(2,1,1,1)/\sqrt7$,
including its signed minimum and even core.
The separate module \code{SharpStabilityExamples.lean} proves the optimal-constant witness
$(a,t,t)$, with $a^2+2t^2=1$ and $0<t\le a/10$, strictly inside the
unique saturated phase. It verifies the actual cap optimum, actual gap,
distance equality, and a finite-dimensional lower bound against every
saturated comparison, including after zero padding. Exact feasible
rational parameters are supplied, so this family is not an empty
conditional implication.

The new modules compile with Lean 4.19.0, \code{--trust=0}, and kernel
type checking enabled. The printed dependency lists for the stated
results contain only the standard logical axioms \code{propext},
\code{Classical.choice}, and \code{Quot.sound}. The accompanying fresh
replay record supplies the exact source hashes, dependencies, and checked
declarations. The nine inherited support-gap modules imported from the
unpublished companion~\cite{Prior} and the $62$ modules of the original
package~\cite{Original} are included and recompiled in the same replay. No claim of a kernel-checked entire manuscript is made.

The separate \code{FourShiftLaw.lean} and \code{FourShiftEnvelope.lean}
derive the actual three-head/four-shift mixture, the strict-interval
identity as open tail minus closed tail, residual normalization and cap,
and the scalar max/min envelope. Their Gaussian comparisons, strict-tail
profile, exponential estimate, above-one tail estimate, and additional
open-tail bounds are explicit named hypotheses. This checks the finite-law
interface conditionally; the full probability optimization and validated
interval certificates remain outside Lean.

\paragraph{Computational scope.}
The finite phase comparisons in Sections~\ref{sec:positive} and
\ref{sec:four-core-full} are reproduced with Python's exact rational
arithmetic. Section~\ref{sec:lower-certificates} gives the full interface
and enclosure algorithm for the validated real interval certificates.
Their acceptance proves the explicit scalar inequalities on complete
rational-box partitions; it does not prove the external Keller--Klein,
Hollom--Portier, or Tomaszewski inputs. These analytic probability
reductions and interval computations are not formalized end to end in
Lean. Floating enumeration and randomized searches were used to falsify
candidate statements during research; they are diagnostics and supply
no proof of the theorems.

\section{Remaining probability questions}\label{sec:limits}
The support-gap problem is fully optimized for every cap, and its
near-maximizers now have the sharp quantitative structure proved above.
The full distributional problem requires an additional step: at a
threshold containing positive central mass, its probability cost must
be compared with the increase in Gaussian mass. No deterministic
support-gap theorem alone supplies that comparison.

The full Kolmogorov optimum is now determined for every
$a\ge1/\sqrt5$. The small-cap theorem in the companion has an
ineffective cutoff, leaving a genuine intermediate range below
$1/\sqrt5$. The balanced matching construction used in the four-core
proof extends to larger even cores, but localization and control of
the remaining core layers become additional obligations. The current
finite certificates establish exactly their stated domains; they are
not evidence of a proved all-cap probability formula.

The exact flat modulus is attained under the stated cap constraints;
at a flat endpoint $a=1/\sqrt k$, its nonflat equality family is
inadmissible. The saturated modulus is exact in the stated local range
$e\le(\sqrt2-1)r$. No assertion is made that these separate curves
already solve the full coefficient-distance optimization at every cap
and every deficit.

\section*{Research process and AI disclosure}

This is fully autonomous agentic research. Dev D. Goyal is responsible
for publication. No expert validation or external peer review is
claimed, and the novelty assessment is provisional.

The results, their proofs, the Lean development, the interval
certificates and the literature comparison were produced by autonomous AI
research agents in an earlier session. Its records describe an autonomous
GPT-6 Astra Pro session with parallel research agents. The proof audits,
counterexample searches and source comparisons recorded in that session
are internal AI-assisted reviews; they are not included here.

Opus 5.5 by Anthropic used with Claude Code assessed the package for this
release. It checked the written proofs, compared the Lean statements with
the theorems, ran fresh Lean replays and axiom audits, reran the exact
rational checks and the interval certificates at $128$ and $192$ bits, and
revised this manuscript: it removed references to files not included here,
stated the dependence on the unpublished companion~\cite{Prior}, and
gathered the disclosure into this section. AI checks are not human peer
review.

\begin{thebibliography}{10}
\interlinepenalty=10000
\bibitem{KK} N. Keller and O. Klein.
\emph{Proof of Tomaszewski's conjecture on randomly signed sums}.
Advances in Mathematics 407 (2022), 108558.
\href{https://arxiv.org/abs/2006.16834}{arXiv:2006.16834}.

\bibitem{HP} L. Hollom and J. Portier.
\emph{Tight anti-concentration of Rademacher sums}.
Random Structures \& Algorithms 67 (2025), e70024.
\href{https://doi.org/10.1002/rsa.70024}{doi:10.1002/rsa.70024}.

\bibitem{DDS} A. De, I. Diakonikolas, and R. A. Servedio.
\emph{A robust Khintchine inequality, and algorithms for computing optimal
constants in Fourier analysis and high-dimensional geometry}.
SIAM Journal on Discrete Mathematics 30(2) (2016), 1058--1094.
\href{https://doi.org/10.1137/130919143}{doi:10.1137/130919143}.

\bibitem{FW} X. Fang and T. Wang.
\emph{A two-regime Khintchine inequality and an improved bound on the
degree-1 Fourier weight for linear threshold functions}.
Preprint, 28 August 2026.
\href{https://arxiv.org/abs/2608.27908}{arXiv:2608.27908}.

\bibitem{Xie} X. Xie.
\emph{Stability of Szarek's inequality with best constant}.
Preprint, 5 October 2026.
\href{https://arxiv.org/abs/2610.07526}{arXiv:2610.07526}.

\bibitem{Pinelis} I. Pinelis.
\emph{On the supremum of the tails of normalized sums of independent
Rademacher random variables}.
Statistics \& Probability Letters 99 (2015), 131--134.
\href{https://arxiv.org/abs/1204.1761}{arXiv:1204.1761}.

\bibitem{Bell} L. Bell.
\emph{Small Lyapunov ratios and the structure of Berry--Esseen maximizers}.
Preprint, 6 October 2026, in the independent Berry--Esseen bounds collection.
\url{https://github.com/epicycloids/independent-berry-esseen-bounds}.
Used for comparison only, not as a proof input.

\bibitem{Ivanisvili} P. Ivanisvili.
\emph{The Keller--Klein Berry--Esseen conjecture: the reachable half for
all Rademacher sums, and the full conjecture for $a_1\ge0.07$}.
Computer-assisted AI-generated manuscript, 21 October 2026.
\url{https://paatai.github.io/ai-proofs/keller-klein/keller-klein.pdf}.

\bibitem{Original} D. D. Goyal.
\emph{The small-coefficient upper half of the Keller--Klein conjecture:
analytic completion and Lean 4 formalization}.
Autonomous AI research package, 30 October 2026.
\url{https://github.com/dvgyl/math-and-theory/tree/main/papers/mathematics/keller-klein}.

\bibitem{Prior} D. D. Goyal.
\emph{Exact extremizers and sharp stability in Rademacher normal approximation}.
Unpublished expanded research package preceding the present extension,
1 November 2026. It is not part of this repository; the nine Lean modules of
it that the present development imports are included with this manuscript.

\bibitem{Arb} F. Johansson.
\emph{Arb: efficient arbitrary-precision midpoint-radius interval arithmetic}.
IEEE Transactions on Computers 66(8) (2017), 1281--1292.
\href{https://doi.org/10.1109/TC.2017.2690633}{doi:10.1109/TC.2017.2690633}.

\bibitem{PythonFlint} The python-flint developers.
\emph{python-flint documentation: real ball arithmetic}.
\url{https://python-flint.readthedocs.io/en/latest/arb.html}.
Version 0.9.0 used for the accompanying interval certificates.
\end{thebibliography}
\end{document}
